% concatenated from main.tex, macros.tex, intro.tex, BiAlgUnivAb.tex, deglocus.tex, Galorbit.tex, PointCount.tex, AppAS.tex, RMM.tex, UnifMM.tex, BettiSubmersive.tex, Specialization.tex, Appendix.tex
\documentclass[12pt,oneside]{amsart}
\usepackage{amsmath,amssymb,cite,mathrsfs,tikz-cd}
\usepackage[all]{xy}
\usepackage{typearea} %%Makes better use of the page, IMHO
\usepackage{hyperref} %%Clickable reference in PDF
\usepackage{color} %% For colorful text
%\usepackage[author={Max Schlepzig}]{pdfcomment}
%\usepackage{refcheck}
%% %%%% Uncomment next three lines for better screen legibility
%% \usepackage{mathpazo}
%% \usepackage{xcolor}
%% \pagecolor[HTML]{2d2d2d}
%% \color[HTML]{cccccc}

\newcounter{maincounter}
\numberwithin{maincounter}{section}
\numberwithin{equation}{section}

\newtheorem{lemma}[maincounter]{Lemma}
\newtheorem{proposition}[maincounter]{Proposition}
\newtheorem{corollary}[maincounter]{Corollary}
\newtheorem{remark}[maincounter]{Remark}
\newtheorem{theorem}[maincounter]{Theorem}
\newtheorem{conjecture}[maincounter]{Conjecture}
\newtheorem{definition}[maincounter]{Definition}



\def\AA{\mathbb{A}}
\def\BB{\mathbb{B}}
\def\EE{\mathbb{E}}
\def\HH{\mathbb{H}}
\def\DD{\mathbb{D}}
\def\NN{\mathbb{N}}
\def\RR{\mathbb{R}}
\def\TT{\mathbb{T}}
\def\CC{\mathbb{C}}
\def\ZZ{\mathbb{Z}}
\def\PP{\mathbb{P}}
\def\QQ{\mathbb{Q}}
\def\FF{\mathbb{F}}
\def\GG{\mathbb{G}}
\def\LL{\mathbb{L}}
\def\MM{\mathbb{M}}
\def\SS{\mathbb{S}}
\def\UU{\mathbb{U}}
\def\XX{\mathbb{X}}

%\def\cal{\mathcal}
\newcommand{\cal}{\mathcal}

\newcommand{\scrD}{\mathscr{D}}
\newcommand{\scrX}{\mathscr{X}}
\newcommand{\scrZ}{\mathscr{Z}}
\newcommand{\cA}{\cal{A}}
\newcommand{\cP}{\cal{P}}
\newcommand{\cU}{\cal{U}}
\newcommand{\cB}{\cal{B}}

%%% Philipp's macros

\newcommand{\dom}[1]{{\mathrm {dom}}({#1})}
\newcommand{\sman}[1]{{#1}^{\mathrm{sm,an}}}
\newcommand{\sm}[1]{{#1}^{\mathrm{sm}}}
\newcommand{\anE}{\mathrm{an}}
\newcommand{\an}[1]{{#1}^{\anE}}
\newcommand{\stab}[1]{{\mathrm{Stab}(#1)}}


\newcommand{\hgtexp}{S}

\newcommand{\rank}{{\rm rank}\,}
\newcommand{\Hpoly}[2]{{H^{}_{#1}({#2})}}
\newcommand{\poly}[2]{{#1^{}({#2})}}
\newcommand{\polyt}[2]{{#1^{\sim}({#2})}}
\newcommand{\polytiso}[2]{{#1^{\sim,{\rm iso}}({#2})}}
%\renewcommand{\graph}[1]{\Gamma({#1})}
\newcommand{\atopx}[2]{{\genfrac{}{}{0pt}{}{#1}{#2}}}
\newcommand{\IP}{{\PP}}
\newcommand{\IG}{{\GG}}
\newcommand{\IC}{{\CC}}
\newcommand{\IR}{{\RR}}
\newcommand{\IF}{{\FF}}
\newcommand{\IT}{{\TT}}
\newcommand{\IRan}{{{\RR}_{\rm an}}}
\newcommand{\IRanexp}{{{\RR}_{\rm an,exp}}}
\newcommand{\RRan}{{\IRan}}
\newcommand{\RRanexp}{{\IRanexp}}
\newcommand{\IRalg}{{\RR}_{\rm alg}}
\newcommand{\IQbar}{{\overline{\QQ}}}
\newcommand{\Kbar}{{\overline{K}}}
\newcommand{\IZ}{{\ZZ}}
\newcommand{\IN}{{\NN}}
\newcommand{\IA}{{\AA}}
\newcommand{\IQ}{{\QQ}}
\newcommand{\IQpbar}{{\overline{\QQ}_p}}
\newcommand{\IQp}{{\QQ_p}}
\newcommand{\ts}[1]{{T}_0({#1})}

\newcommand{\cC}{{\mathcal C}}
\newcommand{\cE}{{\mathcal E}}
\newcommand{\cH}{{\mathcal H}}
\newcommand{\cF}{{\mathcal F}}
\newcommand{\cJ}{{\mathcal J}}
\newcommand{\cK}{{\mathcal K}}
\newcommand{\cL}{{\mathcal L}}
\newcommand{\cM}{{\mathcal M}}
\newcommand{\cO}{{\mathcal O}}
\newcommand{\cV}{{\mathcal V}}
\newcommand{\cW}{{\mathcal W}}
\newcommand{\cX}{{\mathcal{X}}}
\newcommand{\cY}{{\mathcal Y}}
\newcommand{\cZ}{{\mathcal Z}}



\newcommand{\defZ}{Z}
\newcommand{\defF}{F}
\newcommand{\defW}{W}
\newcommand{\defC}{C}
\newcommand{\defE}{E}
%\newcommand{\deffam}{F}


\newcommand{\re}[1]{{\rm Re}({#1})}
%% \newcommand{\imS}{{\rm Im}}
%% \newcommand{\im}[1]{\imS({#1})}
%% \newcommand{\imageS}{{\rm im}}
%% \newcommand{\image}[1]{\imageS({#1})}
\newcommand{\volS}{{\rm vol}}
\newcommand{\vol}[1]{\volS({#1})}
\newcommand{\orth}[1]{{#1}^{\bot}}
\newcommand{\mat}[2]{{\rm Mat}_{#1}({#2})}
\newcommand{\ssm}{\setminus}
\newcommand{\ord}[1]{{\rm ord}({#1})}
\newcommand{\opt}[2]{{\rm Opt}_{#2}({#1})}
\newcommand{\Height}[1]{{H}({#1})}
\newcommand{\trdeg}[1]{{\mathrm{trdeg}}({#1})}
\newcommand{\geo}[1]{\langle {#1}\rangle_{{\rm geo}}}
\newcommand{\defect}{\delta}
\newcommand{\geodef}{{\delta_{\rm geo}}}
\newcommand{\en}[1]{{\rm End}({#1})}
\newcommand{\Hom}[1]{{\rm Hom}({#1})}
\newcommand{\hommaxR}[1]{\text{\rm Hom}({#1})^{*}_{\IR}}
\newcommand{\arith}{\rm arith}
\newcommand{\sgu}[2]{{#1}^{[{#2}]}}
\newcommand{\oa}[1]{{#1}^{\rm oa}}
\newcommand{\codim}{{\rm codim}}
\newcommand{\lgo}{LGO}
\newcommand{\zcl}[1]{{\rm Zcl}({#1})}


\newcommand{\trans}[1]{{#1}^{T}}

\newcommand{\red}[1]{\textcolor{red}{#1}}

\renewcommand{\subset}{\subseteq} %%% Some people think \subset
%%% excludes equality
\renewcommand{\supset}{\supseteq}

\newcommand{\gra}[1]{\mathrm{Gr}({#1})}


\newcommand{\gl}[2]{{\mathrm {GL}}_{#1}({#2})}
\renewcommand{\sp}[2]{{\mathrm {Sp}}_{#1}({#2})}
\newcommand{\autS}{{\mathrm {Aut}}}
\newcommand{\aut}[1]{\autS({#1})}

\newcommand{\spec}[1]{\mathrm{Spec}\,{#1}}

\newcommand{\tor}[1]{{#1}_{\mathrm{tor}}}
\newcommand{\gal}[1]{{\mathrm{Gal}}({#1})}

\newcommand{\zeroset}[1]{\mathscr{Z}({#1})}

\newcommand{\jac}[1]{\mathrm{Jac}({#1})}
\newcommand{\diffS}{\mathrm{Diff}}

\newcommand{\diffn}[1]{\diffS_n({#1})}

\newcommand{\cg}[1]{\mathrm{Cl}({#1})}

\newcommand{\maxtr}{\mathrm{maxtr}\,}
\setcounter{tocdepth}{1}

\newcommand{\tr}[2]{{\mathrm{Tr}}_{{#1}}({#2})}
\newcommand{\im}[2]{{\mathrm{Im}}_{{#1}}({#2})}

\newcommand{\pullback}[1][dr]{\save*!/#1-1.2pc/#1:(-1,1)@^{|-}\restore}
\newcommand{\bigpullback}[1][dr]{\save*!/#1-2.2pc/#1:(-2,2)@^{|-}\restore}
\newcommand{\leftpullbak}[1][dl]{\save*!/#1-1.2pc/#1:(-1,1)@^{|-}\restore}

%\newcommand{\rk}[1]{\mathrm{rk}({#1})}

\newcommand{\ntconst}{c_\mathrm{NT}}

\newcommand{\pdfhl}[2]{\pdfmarkupcomment[markup=Highlight,color=yellow]{{#1}}{{#2}}}

\newcommand{\pullbackcorner}[1][dr]{\save*!/#1-1.7pc/#1:(-1.5,1.5)@^{|-}\restore}

\newcommand{\philipp}[1]{{\color{blue} Philipp: #1}}
\newcommand{\ziyang}[1]{{\color{red} Ziyang: #1}}

\DeclareMathOperator{\rk}{rk}

%\makeindex  %%% Generates the index. Remove for final version.

\begin{document}
\title{The relative Manin--Mumford conjecture}
\author{Ziyang Gao}
\author{Philipp Habegger}

\address{Department of Mathematics, UCLA, Los Angeles, CA 90095, USA}
\email{ziyang.gao@math.ucla.edu}
\address{Department of Mathematics and Computer Science; University of Basel; Spiegelgasse 1, 4051 Basel, Switzerland}
\email{philipp.habegger@unibas.ch}


\subjclass[2000]{11G30, 11G50, 14G05, 14G25}

\maketitle

{\centering\footnotesize \textit{To the memory of Bas Edixhoven}\par}

\begin{abstract}
  We prove the Relative Manin--Mumford Conjecture for
  families of abelian varieties in characteristic $0$. We follow the
  Pila--Zannier method to study special point problems, and we use the Betti map which 
  goes back to work of Masser and Zannier in the case of curves. The key
  new ingredients compared to previous applications of this approach are
  a height inequality proved by both authors of the current paper and
  Dimitrov, and the first-named author's study of certain degeneracy
  loci in subvarieties of abelian schemes.
  
  We also strengthen this result and prove a criterion for torsion points to be dense in a subvariety of an abelian scheme over $\IC$.

  The Uniform Manin--Mumford
  Conjecture for curves embedded in their Jacobians was first proved
  by K\"uhne. We give a new  proof, as a corollary to our main theorem, that does
  not use equidistribution.  
% We also strengthen this result and prove a criterion for torsion points to be dense in a subvariety of an abelian scheme.%, and 
\end{abstract}
\tableofcontents


%% Section 1
\section{Introduction}

Let $S$ be a regular, irreducible, quasi-projective
variety defined over an algebraically closed  field $L$ of characteristic $0$. Let $\pi
\colon \cA \rightarrow S$ be an abelian scheme of relative dimension
$g  \ge 1$, namely a proper smooth group scheme whose fibers are
abelian varieties.
Let $\cA_{\mathrm{tor}}$ denote the union over all $s\in S(L)$ of
the group of all torsion points in $\cA_s = \pi^{-1}(s)$. For each $N
\in \mathbb{Z}$, let $[N] \colon \cA \rightarrow \cA$ be the
multiplication-by-$N$ morphism.

The goal of this paper is to prove the \textit{relative Manin--Mumford
  conjecture} for abelian schemes.

\begin{theorem}
  \label{MainThm}
  Let $X$ be an irreducible subvariety of $\cA$. Assume that
  % $X_{\overline{\eta}}$ is irreducible
  % each irreducible
  % component of
  % and not contained in any proper
  % algebraic subgroup of $\cA_{\overline{\eta}}$. 
  $\mathbb{Z}X := \bigcup_{N\in \mathbb{Z}} [N]X$ is Zariski dense in
  $\cA$.
  If $X(L) \cap \cA_{\mathrm{tor}}$ is Zariski dense in $X$, then
  $\dim X \ge g$.
\end{theorem}
Throughout the whole paper, by \textit{irreducible subvariety} we mean \textit{closed irreducible subvariety} unless stated otherwise.


The assumption $\ZZ X$ being Zariski dense in $\cA$ can be checked
over the geometric generic fiber of $\cA \rightarrow S$. Indeed, let
$\eta$ be the generic point of $S$ and fix an algebraic closure of the
function field of $S$. Write $X_{\overline{\eta}}$ for the geometric
generic fiber of $\pi|_X$.
Then $X_{\overline \eta}$ is non-empty if and only if $\pi|_X\colon
X\rightarrow S$ is dominant.
In particular, $\cA_{\overline{\eta}}$ is an abelian variety over
an algebraically closed field containing the possible reducible $X_{\overline\eta}$. Then $\ZZ X$ is Zariski dense in $\cA$
if and only if $X_{\overline{\eta}}$ is non-empty and not contained in
a finite union of proper algebraic subgroups of $\cA_{\overline{\eta}}$.


The Relative Manin--Mumford Conjecture was inspired by S.~Zhang's ICM
talk \cite{zhang1998small} and proposed by Pink \cite[Conjecture~6.2]{Pink} and %Zannier 
%\cite[Conjecture~6.2]{Pink}. The current formulation was proposed by
Zannier~\cite{ZannierBook}. In the case $\dim X=1$ it was proved in a
series of papers by Masser--Zannier and Corvaja--Masser--Zannier
\cite{MZ:torsionanomalous,MasserZannierTorsionPointOnSqEC,
MASSER2014116, MasserZannierRelMMSimpleSur, CorvajaMasserZannier2018,
MasserZannierRMMoverCurve}. See also work of
Stoll~\cite{Stoll:simtorsion} for an explicit case. For surfaces some 
results are due to the first-named author~\cite{hab:weierstrass} and the recent work of 
Corvaja--Tsimerman--Zannier~\cite{CTZ:23}. When $\cA$ is a fibered product of families of elliptic curves, it was proved by K\"{u}hne \cite{KuehneRBC}. A core idea of many of these papers, as in ours, is to 
use the Betti coordinates introduced by Masser and Zannier~\cite{MZ:torsionanomalous} to study questions in diophantine
geometry.
%The core method of many of these papers, as in ours, is ultimately based on the approach
%initiated by Masser and Zannier~\cite{MZ:torsionanomalous} who
%introduced the use of Betti coordinates to questions in diophantine
%geometry.

As a corollary  of Theorem~\ref{MainThm} we obtain  the \textit{Uniform
Manin--Mumford Conjecture} for curves embedded in their Jacobians, which was recently proved by K\"{u}hne
\cite[Theorem~1.2]{KuehneUnifMM}.

\begin{corollary}\label{CorUniformMM}
  For each integer $g \ge 2$, there exists a constant $c = c(g) > 0$ with the following property. 
  Let $C$ be an irreducible, smooth, projective curve of genus $g$
  defined over $\IC$. Let $x_0 \in C(\IC)$, and let $C-x_0$ be the
  image of the Abel--Jacobi embedding based at $x_0$ in the Jacobian
  $\mathrm{Jac}(C)$ of $C$. Then
  \begin{equation}
    \#(C(\IC)-x_0) \cap \mathrm{Jac}(C)_{\mathrm{tor}} \le c.
  \end{equation}
\end{corollary}

%Corollary~\ref{CorUniformMM} was recently proved by K\"{u}hne
%\cite[Theorem~1.2]{KuehneUnifMM}. 
A second proof of the Uniform Manin--Mumford Conjecture for curves was given by Yuan in \cite{YuanArithBig}, based on the theory of adelic line bundles over 
quasi-projective varieties of Yuan--Zhang \cite{YuanZhangEqui}. Prior to K\"{u}hne's 
proof of the full conjecture, DeMarco--Krieger--Ye \cite{DeMarcoKriegerYeUniManinMumford} proved the case where $g = 2$ and $C$ is bi-elliptic, using method of 
arithmetic dynamical systems.


We hereby give a different proof. The
common tools used in \cite{KuehneUnifMM} and the current paper are the
height inequality \cite[Theorems 1.6 and B.1]{DGHUnifML} proved by both
authors of the current paper and Dimitrov, and the first-named
author's results on the generic rank of the Betti map \cite{GaoBettiRank}. K\"{u}hne used this height inequality,
among other tools, to prove an equidistribution result. 
He then applied this equidistribution result and %the first-named author's
\cite[Theorem 1.3]{GaoBettiRank} in combination with the Ullmo--S.~Zhang
approach to the Bogomolov Conjecture to conclude the argument.
% the Uniform Manin--Mumford Conjecture.
% on an appropriate abelian scheme.

Our proof does not involve equidistribution, and we apply
\cite{GaoBettiRank} in a different way.  The crucial case for
Theorem~\ref{MainThm} is the case where $L$ is an algebraic closure
$\IQbar$ of $\IQ$, \textit{i.e.}, in the setup of the theorem we assume every variety is defined over $\IQbar$. The proof of Theorem~\ref{MainThm} over $\IQbar$ occupies the current paper up to $\mathsection$\ref{SectionConclusion}. Then in
$\mathsection$\ref{SectionUnifMM} we prove
Corollary~\ref{CorUniformMM} as a consequence of
Theorem~\ref{MainThm} over $\IQbar$. The deduction is inspired by \cite[Theorem~2.4]{Stoll:Uniform} and \cite[Proposition~7.1]{DGHUnifML}.

%<<<<<<< HEAD
Our proof of Theorem~\ref{MainThm} for  $L=\IQbar$ is in spirit of the Pila--Zannier
% =======
% Our proof of Theorem~\ref{MainThm} is in spirit of the Pila--Zannier
% >>>>>>> parent of fddc4d3 (Added a section for the criterion in terms of Betti rank. Added comments in Introduction)
method~\cite{PilaZannier} to solve special point problems. Roughly
speaking this strategy can be divided into four steps. We start  
%by deducing 
with a 
large Galois orbit result on 
%from work of 
%David~\cite{DavidMinHaut}
%R\'{e}mond } on 
%points of small N\'eron--Tate height 
torsion points on an abelian variety,
quantifying earlier work of Masser~\cite{Masser:smallvalues}. Such a
result can be deduced from work of David~\cite{DavidMinHaut}. It
also follows 
 more directly by R\'{e}mond and Gaudron's refinement
 \cite{RemondGaloisBound, GRGaloisBound}  of deep work of Masser and
 W\"{u}stholz
 on isogeny estimates for abelian varieties \cite{MW:abelianisog}. 
%A weaker result in this direction which suffices for our purpose 
%is proved in the Appendix using a result of David~\cite{DavidMinHaut}. 
A key new input at this step is the 
height
inequality \cite[Theorem~B.1]{DGHUnifML} which roughly speaking allows
us to bound the height of the abelian variety itself.
Then we introduce a suitable set that is definable in an o-minimal
structure, here $\IRanexp$, that encodes our points of interest as rational.
We then use this result to invoke an
appropriate version of the Pila--Wilkie counting theorem
due to the second-named author and Pila \cite[Corollary~7.2]{HabeggerPilaENS}.
Finally, we apply a suitable functional transcendence theorem, the
mixed Ax--Schanuel theorem proved by the first-named author
\cite{GaoMixedAS}. In this last step we study the degeneracy locus,
defined in $\mathsection$\ref{secdeglocus}, as was done in
\cite[Proposition~1.10]{GaoBettiRank}.
 
Our application of the height
inequality \cite[Theorem~B.1]{DGHUnifML} differs  from the approaches in
\cite{DGHUnifML} and \cite{KuehneUnifMM}. Instead of constructing a
non-degenerate subvariety, we study the degeneracy loci more
carefully. We prove the desired result by dividing into two cases:
either the height inequality \cite[Theorem~B.1]{DGHUnifML} is
applicable or not. If it is not applicable, then the $0$-th degeneracy
locus is large and we may apply \cite[Proposition~1.10]{GaoBettiRank}.
If it is applicable, then we follow the Pila--Zannier method described
above. Ultimately we show that 
\cite[Proposition~1.10]{GaoBettiRank} still applies and can conclude
the Relative Manin--Mumford Conjecture.


\subsection{Criterion of torsion points being dense}
%Let $\pi \colon \cA \rightarrow S$ be an abelian scheme defined over $\IC$ as at the beginning of this paper, and let $X$ be an irreducible subvariety of $\cA$.
Let us assume for the moment that the base field $L$ is a subfield of $\IC$.

It is natural to ask whether the following converse of
Theorem~\ref{MainThm} is true.
Let $X\subset \cA$ be an irreducible subvariety.
If $\ZZ X$ is Zariski
dense in $\cA$ and $\dim X \ge g$, then is it true that $X(\IC) \cap
\cA_{\mathrm{tor}}$ is Zariski dense in $X$? This question is
related to the generic Betti rank, as is studied in \cite{ACZBetti}.
Indeed, using \cite[Proposition~2.1.1]{ACZBetti}, one can show that the
answer is yes \textit{in some cases}, even for the Euclidean topology. For example three cases were proved in \cite{ACZBetti}: if $g = 2, 3$ and $\cA/S$ has no fixed part over any finite covering of $S$, if 
the geometric generic fiber $\cA_{\overline{\eta}}$ satisfies 
$\mathrm{End}(\cA_{\overline{\eta}}) = \ZZ$, and if $\cA \rightarrow
S$ is the Jacobian of the universal hyperelliptic curve.
\cite[Theorem~1.4.(i)]{GaoBettiRank} proves some more general cases,
for example if $\cA_{\overline{\eta}}$ is simple. Some new cases for
the denseness in the Euclidean topology were
recently proved by Eterovi\'{c}--Scanlon in \cite{EterovicScanlon}.% in.
%!!!EUCLIDEAN TOPOLOGY IN [ES].

In this paper, we provide a criterion for $X(\IC) \cap \cA_{\mathrm{tor}}$ being Zariski dense in $X$. %
Using this criterion, one can show that the converse of
Theorem~\ref{MainThm} is \textit{false in general} even if $\cA/S$ has
no fixed part over any finite covering of $S$. See  \cite[Example~9.4]{GaoBettiRank} for a counterexample with $g = 4$. 
 The statement of this criterion, Theorem~\ref{ThmCritTorsionDense}, involves the Betti map and the Betti
 rank. We briefly recall the definition here. For a precise definition of the Betti map, we refer to \cite[$\mathsection$3-$\mathsection$4]{GaoBettiRank} or \cite[$\mathsection$2.3 and $\mathsection$B.1]{DGHUnifML}.

For any $s \in S(\IC)$, there exists a simply-connected open neighborhood $\Delta
\subseteq S^{\mathrm{an}}$ of $s$; here and below the superscript
``$\mathrm{an}$'' refers to complex analytification. Then one can define a basis $\omega_1(s),\ldots,\omega_{2g}(s)$ of the period lattice of each fiber $s \in \Delta$ as holomorphic functions of $s$. Now each fiber $\cA_s = \pi^{-1}(s)$ can be identified with the complex torus $\IC^g/\left(\IZ \omega_1(s)\oplus \cdots \oplus \IZ\omega_{2g}(s)\right)$, and each point $x \in \cA_s(\IC)$ can be expressed as the class of $\sum_{i=1}^{2g}b_i(x) \omega_i(s)$ for real numbers $b_1(x),\ldots,b_{2g}(x)$. Then $b_{\Delta}(x)$ is defined to be the class of the $2g$-tuple $(b_1(x),\ldots,b_{2g}(x)) \in \IR^{2g}$ modulo $\IZ^{2g}$.
We obtain a real-analytic map $
b_{\Delta} \colon \cA_{\Delta} = \pi^{-1}(\Delta) \rightarrow \mathbb{T}^{2g}$, 
which is fiberwise a group isomorphism and
where $\mathbb{T}^{2g}$ is the real torus of dimension $2g$. The map
$b_{\Delta}$ is well-defined up-to real analytic automorphisms of
$\mathbb{T}^{2g}$, \textit{i.e.}, elements of $\mathrm{GL}_{2g}(\ZZ)$.

Let $X^{\mathrm{reg}}$ denote the regular locus of $X$.
The \textit{generic Betti rank} of $X$ is defined in terms of the
differential
$\mathrm{d} b_\Delta$ to be
\[
\mathrm{rank}_{\mathrm{Betti}}(X) := \max_{x \in X^{\mathrm{reg,an}}\cap \cA_{\Delta}} \mathrm{rank}_{\RR} (\mathrm{d}b_{\Delta}|_{X\cap \cA_{\Delta}})_x.
\]
This definition does not depend on the choice of $\Delta$ (see \cite[end of $\mathsection$4]{GaoBettiRank}) or $b_{\Delta}$ (since every two $b_{\Delta}$ differ from an element in $\mathrm{GL}_{2g}(\ZZ)$). The definition of the generic Betti rank yields the following trivial upper bound
\begin{equation}\label{EqBettiRank}
\mathrm{rank}_{\mathrm{Betti}}(X) \le 2 \min\{\dim X, g\}.
\end{equation}

%!!!I CHANGED TO THEOREM, $\IC$ IS NOW $L$.
\begin{theorem}
  \label{ThmCritTorsionDense}
  Assume $\ZZ X$ is Zariski dense in $\cA$. Then $X(L) \cap \cA_{\mathrm{tor}}$ is Zariski dense in $X$ if and only if $\mathrm{rank}_{\mathrm{Betti}}(X) = 2g$.
\end{theorem}

The ``if'' direction of this theorem is proved by
\cite[Proposition~2.1.1]{ACZBetti} and uses real analytic geometry. In fact, if $\mathrm{rank}_{\mathrm{Betti}}(X) = 2g$, then $X(\IC) \cap \cA_{\mathrm{tor}}$ is dense in the Euclidean topology on $X^{\mathrm{an}}$.

The ``only if'' direction of this theorem implies Theorem~\ref{MainThm} by \eqref{EqBettiRank}. In this paper, we prove this direction by a combination of Theorem~\ref{MainThm} and the criterion for $\mathrm{rank}_{\mathrm{Betti}}(X) = 2g$ given by \cite[Theorem~1.1]{GaoBettiRank}.


\subsection{From $\IQbar$ to $\IC$}
Early work of Bombieri--Masser--Zannier \cite{BMZUnlikely} shows how
to reduce certain Unlikely Intersection statements from base field $\IC$ to
$\IQbar$ using a specialization argument.
Barroero and Dill~\cite{BarDill:distinguised} provide a
 framework to approach these question in large generality. 
A specialization argument, based on a result of Masser, allows to pass
from $\IQbar$ to $\IC$ for Corollary~\ref{CorUniformMM}. However, the
specialization argument for Theorem~\ref{MainThm} and
Theorem~\ref{ThmCritTorsionDense} is more complicated.

All the varieties in question are defined over an algebraically closed
field of finite transcendence degree $d$ over $\IQbar$.
The proof of Theorem~\ref{MainThm} for $L=\IQbar$, \textit{i.e.}, if $d =
0$ is completed in $\mathsection$\ref{SectionConclusion}.
The proof of Theorem~\ref{ThmCritTorsionDense} is completed in
$\mathsection$\ref{SectionCriterionTorsionDenseOverIQbar}.


After this, we proceed by induction on $d$ by proving 
  two things in
$\mathsection$\ref{SectionSpecialization}: Theorem~\ref{MainThm} when
$\mathrm{trdeg}_\IQbar L\le d$ implies
Theorem~\ref{ThmCritTorsionDense} when
$\mathrm{trdeg}_\IQbar L\le d$, and
Theorem~\ref{ThmCritTorsionDense} when
$\mathrm{trdeg}_\IQbar L\le d$ implies
Theorem~\ref{MainThm} for
$\mathrm{trdeg}_\IQbar L\le d+1$.


Shortly before this paper was finalized, Corvaja, Tsimerman, and Zannier~\cite[Appendix~A]{CTZ:23} have independently
shown how to reduce Theorem~\ref{MainThm} from
$\IC$ to $\IQbar$ with a different argument.




%, and this implies the full Theorem~\ref{MainThm} by \eqref{EqBettiRank}.

\subsection{Notation} 
Let $\mathbb{A}_g$ be the moduli space of abelian varieties of
dimension $g$ with a polarization of type $\mathrm{diag}(d_1,\ldots,d_g)$
with $d_1\mid \cdots\mid d_g$ positive integers, endowed with
symplectic level-$\ell$-structure for some large enough but fixed
$\ell$ that is coprime to $d_g$ (so that $\mathbb{A}_g$ is a fine
moduli space).
Let $\pi\colon \mathfrak{A}_g \rightarrow \mathbb{A}_g$
be the universal abelian variety.
Both $\mathfrak{A}_g$ and $\IA_g$ are irreducible, regular,
quasi-projective varieties definable over a number field.

Any abelian variety is isogenous to a principally polarized abelian
variety, after extending the base field. We may apply this observation
to the generic fiber of our abelian schemes. Our main results
Theorems~\ref{MainThm} and \ref{ThmCritTorsionDense} are insensitive
to \'etale base change of $S$ and up-to isogeny. Therefore it is
possible to recover all results here while assuming $d_1=\cdots=d_g=1$
for the universal family.

%For $N\in \IZ$ and any abelian scheme $\cA/S$, denote by $[N] \colon \cA \rightarrow \cA$ the multiplication-by-$N$ morphism.

% Let $(\mathfrak{A}_g)_{\mathrm{tor}}$ be the set $x \in \mathfrak{A}_g(\IQbar)$ such that $[N]x$ lies in the zero section of $\mathfrak{A}_g \rightarrow \mathbb{A}_g$.


\subsection*{Acknowledgements} Ziyang Gao has received
funding from the European Research Council (ERC) under the
European Union's Horizon 2020 research and innovation programme (grant
agreement n$^\circ$ 945714). Philipp Habegger has received funding
from the Swiss National Science Foundation (grant n$^\circ$ 200020\_184623). 
The authors would like to thank Daniel Bertrand and Umberto Zannier for having communicated with us the current formulation of the Relative Manin--Mumford Conjecture, in particular during the conference ``Third ERC Research Period on Diophantine Geometry'' in Rome 2015, and would like to thank the organizers of this conference. We would also like to thank Michael Stoll for comments on a previous version of the paper and Eric Gaudron for having pointed out to us the reference \cite{RemondGaloisBound}. We would like to thank the referees for their careful reading and valuable comments.

%%% Local Variables:
%%% TeX-master: "main"
%%% End:


%% Section 2 (new, was 4)
\section{Bi-algebraic Structure on the Universal Abelian Variety}
\label{SectionBiAlg}
The goal of this section is to collect some basic facts about the universal abelian variety, especially its bi-algebraic structure. The end of this section contains the functional transcendence statement, the \textit{Ax--Schanuel theorem} for the universal abelian variety.

Let $d_1,\ldots,d_g$ be positive integers such that $d_1\mid \cdots
\mid d_g$. Let $D = \mathrm{diag}(d_1,\ldots,d_g)$ be the $g\times g$
diagonal matrix.

\subsection{Moduli space of abelian varieties} 
Let $\mathbb{A}_g$ be the moduli space of abelian varieties which are polarized of type $D$. In this subsection, we describe the bi-algebraic geometry on $\mathbb{A}_g$.

Let $\mathfrak{H}_g$ be the Siegel upper half space defined by
\[
\{ Z = X + \sqrt{-1}Y \in \mathrm{Mat}_{g \times g}(\mathbb{C}) : Z =
Z^{\!^{\intercal}},X,Y\in \mathrm{Mat}_{g\times g}(\mathbb{R}),~ Y~\text{is positive definite}\}.
\]
On identifying a complex number with its real and imaginary parts, 
 $\mathfrak{H}_g$ becomes a semi-algebraic subset of $\IR^{2g^2}$.
Note that $\mathfrak{H}_g$ is an open subset, in the Euclidean topology,  of
$M_{g\times g}(\IC) = \{Z \in \mathrm{Mat}_{g \times g}(\mathbb{C}) : Z =
Z^{\!^{\intercal}}\} \cong \IC^{g(g+1)/2}$. In fact,
$\mathfrak{H}_g$ is a connected 
complex manifold. 
It is well-known that the universal covering space of
$\mathbb{A}_g^{\mathrm{an}}$, the analytification of $\mathbb{A}_g$,
is given by $\mathfrak{H}_g$ and the covering map is a holomorphic map
\[
u_B \colon \mathfrak{H}_g \rightarrow  \mathbb{A}_g^{\mathrm{an}}.
\]

% Thus we endow $\mathfrak{H}_g$ with the following algebraic structure, and therefore define a bi-algebraic structure on $\mathbb{A}_g$.

\begin{definition}%\label{DefnBiAlgModuliSpace}
\begin{enumerate}
\item[(i)] A subset $\tilde{Y} \subseteq \mathfrak{H}_g$ is said to be
  irreducible algebraic if $\tilde{Y}$ is a complex  analytic
  irreducible component of $W \cap \mathfrak{H}_g$, where $W$ is an
  algebraic subset of $M_{g \times g}(\mathbb{C})$.
% PH: Never used
  % \item[(ii)] An irreducible algebraic subset $\tilde{Y} \subseteq \mathfrak{H}_g$ is said to be bi-algebraic if $u_B(\tilde{Y})$ is an algebraic subvariety of $\mathbb{A}_g$.
\item[(ii)] An irreducible subvariety $Y$ of $\mathbb{A}_g$ is said to be bi-algebraic if $Y = u_B(\tilde{Y})$ for some irreducible algebraic subset $\tilde{Y}$ of $\mathfrak{H}_g$.
\end{enumerate}
\end{definition}

The following notation will be used. Let $\tilde{Z}$ be a complex
analytic irreducible subset of $\mathfrak{H}_g$. We use
$\tilde{Z}^{\mathrm{Zar}}$ to denote the smallest algebraic subset of
$\mathfrak{H}_g$ that contains $\tilde{Z}$
and use $u_B(\tilde{Z})^{\mathrm{biZar}}$ to denote the smallest
bi-algebraic subvariety of $\mathbb{A}_g$ that contains
$u_B(\tilde{Z})$. We have the following
relation
\begin{equation*}
  %\label{EqBiAlgRelationModuliSpace}
%\tilde{Z}^{\mathrm{Zar}} \subseteq \tilde{Z}^{\mathrm{biZar}} \quad\text{and}\quad
u_B(\tilde{Z})^{\mathrm{Zar}} \subseteq u_B(\tilde{Z})^{\mathrm{biZar}}.
%, \quad u(\tilde{Z}^{\mathrm{biZar}}) = u(\tilde{Z})^{\mathrm{biZar}}.
\end{equation*}




\subsection{Universal abelian variety}
%\label{SubsectionUnivAbVar}
Let $\pi \colon \mathfrak{A}_g \rightarrow \mathbb{A}_g$ be the universal abelian variety.

The uniformization of $\mathfrak{A}_g$, in the category of complex
spaces, is given by theta functions
\begin{equation}\label{EqUnifUniversalAbVar}
u \colon \IC^g \times  \mathfrak{H}_g \rightarrow \mathfrak{A}_g^{\mathrm{an}}.
\end{equation}

As for $\mathfrak{H}_g$, the complex space $\IC^g \times  \mathfrak{H}_g$ is naturally an open subset of the analytification of the complex algebraic variety $\IC^g \times M_{g \times g}(\mathbb{C})$. Thus we endow $\IC^g \times  \mathfrak{H}_g$ with the following algebraic structure, and therefore define a bi-algebraic structure on $\mathfrak{A}_g$.
\begin{definition}%\label{DefnBiAlgUnivAb}
\begin{enumerate}
\item[(i)] A subset $\tilde{Y} \subseteq \IC^g \times  \mathfrak{H}_g$
  is said to be irreducible algebraic if $\tilde{Y}$ is a complex
  analytic irreducible component of $W \cap (\IC^g \times
  \mathfrak{H}_g)$, where $W$ is an algebraic subset of $\IC^g
  \times M_{g \times g}(\mathbb{C})$.
  % %Never used
% \item[(ii)] An irreducible algebraic subset $\tilde{Y} \subseteq \IC^g
% \times \mathfrak{H}_g$ is said to be bi-algebraic if $u(\tilde{Y})$ is
% an algebraic subvariety of $\mathfrak{A}_g$.
\item[(ii)] An irreducible subvariety $Y$ of $\mathfrak{A}_g$ is said to be bi-algebraic if $Y = u(\tilde{Y})$ for some irreducible algebraic subset $\tilde{Y}$ of $\IC^g \times  \mathfrak{H}_g$.
\end{enumerate}
\end{definition}

We refer to \cite[Proposition~5.3]{GaoBettiRank} for a geometric description of bi-algebraic subvarieties of $\mathfrak{A}_g$. 
As for the moduli space, the following notation will be used. Let
$\tilde{Z}$ be a complex analytic irreducible subset of $\IC^g \times
\mathfrak{H}_g$. We use $\tilde{Z}^{\mathrm{Zar}}$ to denote the
smallest algebraic subset of $\IC^g \times \mathfrak{H}_g$ that
contains $\tilde{Z}$
% , use $\tilde{Z}^{\mathrm{biZar}}$ to denote the
% smallest bi-algebraic subset of $\IC^g \times \mathfrak{H}_g$ which
% contains $\tilde{Z}$,
and use $u(\tilde{Z})^{\mathrm{biZar}}$ to
denote the smallest bi-algebraic subvariety of $\mathfrak{A}_g$ that
contains $u(\tilde{Z})$. We have the following relation %The following relations are easy to verify:
\begin{equation}
%  \tilde{Z}^{\mathrm{Zar}} \subseteq \tilde{Z}^{\mathrm{biZar}}, \quad
  u(\tilde{Z})^{\mathrm{Zar}} \subseteq u(\tilde{Z})^{\mathrm{biZar}}.
%  \quad u(\tilde{Z}^{\mathrm{biZar}}) = u(\tilde{Z})^{\mathrm{biZar}}.
\end{equation}



The following lemma summarizes the bi-algebraic structures discussed above.
\begin{lemma}
We have the following  commutative diagram:
\begin{equation*}
  %\label{EqDiagUnifProj1}
\xymatrix{
 \IC^g \times \mathfrak{H}_g  \ar[r]^-{u} \ar[d]_{\tilde{\pi}} & \mathfrak{A}_g^{\mathrm{an}} \ar[d]^{\pi} \\ \mathfrak{H}_g \ar[r]^-{u_B} & \mathbb{A}_g^{\mathrm{an}},
}
\end{equation*}
where $\tilde{\pi}$ is the natural projection (hence is the
restriction of an algebraic map), and both $u$ and $u_B$ are uniformizations in the category of complex spaces.
\end{lemma}


\subsection{Functional transcendence for $\mathfrak{A}_g$}
The following theorem, known as the \textit{weak mixed Ax--Schanuel
  theorem for $\mathfrak{A}_g$}, was
proved by the first-named author.
\begin{theorem}[{\!\!\cite[Theorem~3.5]{GaoMixedAS}}]\label{ThmAS}
Let $\tilde{Z}$ be an irreducible complex analytic 
subset of $\IC^g \times\mathfrak{H}_g$. Then we have
\[
\dim \tilde{Z}^{\mathrm{Zar}} + \dim u(\tilde{Z})^{\mathrm{Zar}} \ge \dim \tilde{Z} + \dim u(\tilde{Z})^{\mathrm{biZar}}.
\]
\end{theorem}
%Let us explain some notation in this theorem: $\tilde{Z}^{\mathrm{Zar}}$ is defined to be the smallest algebraic subset of $\IC^g \times\mathfrak{H}_g$ which contains $\tilde{Z}$ for the algebraic structure defined in Definition~\ref{DefnBiAlgUnivAb}.(i), and $u(\tilde{Z})^{\mathrm{biZar}}$ is defined to be the smallest bi-algebraic subvariety of $\mathfrak{A}_g$ which contains $u(\tilde{Z})$.


\subsection{The universal uniformized Betti map}
The following Betti map has been proved to be important in several Diophantine problems.

Consider the real-algebraic isomorphism
\begin{equation}\label{EqComplexStrOfX2g}
\begin{array}{cccc}
 \mathbb{R}^g \times \mathbb{R}^g \times \mathfrak{H}_g & \xrightarrow{\sim} & \mathbb{C}^g \times \mathfrak{H}_g, \\
 (a,b,Z) & \mapsto & (Da+Zb, Z)
\end{array}.
\end{equation}
The \textit{universal uniformized Betti map} is the semi-algebraic map
\begin{equation}\label{EqUnivBettiMap}
\tilde{b} \colon \IC^{g} \times \mathfrak{H}_g \rightarrow \IR^{2g},
\end{equation}
which is the inverse of \eqref{EqComplexStrOfX2g} composed with the natural projection $ \mathbb{R}^g \times \mathbb{R}^g \times \mathfrak{H}_g \rightarrow \IR^{2g}$.

%%% Local Variables:
%%% TeX-master: "main"
%%% End:


%% Section 3 (new, was 2)
\section{The Degeneracy Locus}
\label{secdeglocus}
\subsection{Degeneracy loci}
An important notion in our approach to prove the Relative
Manin--Mumford Conjecture is the \textit{$t$-th degeneracy locus}
introduced by the first-named author in \cite[Definition~1.6]{GaoBettiRank}.
In this paper, we need the cases $t = 0$ and $t = 1$.

Let $g\ge 1$. We consider $\mathfrak{A}_g$ as a smooth, irreducible,
quasi-projective variety defined over $\IC$. Recall that $\pi\colon
\mathfrak{A}_g\rightarrow\IA_g$ is the structural morphism.

\begin{definition}
  %\label{DefnDegeneracyLocus}
Let $X$ be an irreducible subvariety of $\mathfrak{A}_g$. Let $t \in \IZ$. The $t$-th degeneracy locus, denoted by $X^{\mathrm{deg}}(t)$, is the union
\begin{equation}\label{EqDefnDegLocus}
X^{\mathrm{deg}}(t) = \bigcup_{\substack{Y \subseteq X\text{ irreducible,}~ \dim Y > 0 \\ \dim Y^{\mathrm{biZar}} - \dim \pi(Y)^{\mathrm{biZar}} < \dim Y + t}} Y.
\end{equation}
\end{definition}
We refer to $\mathsection$\ref{SectionBiAlg} for the notations
$Y^{\mathrm{biZar}}$ and $\pi(Y)^{\mathrm{biZar}}$. However, we will
not need this until $\mathsection$\ref{SectionAppAS}. Notice that
\eqref{EqDefnDegLocus} is not precisely \cite[Definition~1.6]{GaoBettiRank}, but these two definitions are equivalent for $X
\subseteq \mathfrak{A}_g$ by \cite[Corollary~5.4]{GaoBettiRank}. It is proved in \cite{GaoBettiRank} that $X$ is non-degenerate, \textit{i.e.} the generic Betti rank of $X$ is $2\dim X$, if and only if $X^{\mathrm{deg}}(0) \not= X$.

An immediate consequence of the definition is $X^{\mathrm{deg}}(0) \subseteq X^{\mathrm{deg}}(1)$.

The degeneracy locus is a possibly infinite union of Zariski closed
subsets. Nevertheless we have the following theorem.
\begin{theorem}[{\!\!\cite[Theorem~1.8]{GaoBettiRank}}]
   \label{ThmDegLocusZarClosed}
   The degeneracy locus $X^{\mathrm{deg}}(t)$ is  Zariski closed
   in $X$ for all $t$.
\end{theorem} 

In the current paper, we apply this theorem in
$\mathsection$\ref{SectionConclusion}. Moreover, if $X$ is defined
over $\IQbar$, then $X^{\mathrm{deg}}(t)$ is defined over $\IQbar$ by
\cite[Theorem~4.2.4]{GaoHDR}; however we will not use this fact in the
current work. %!!!CHECK IF LAST CLAUSE IS TRUE



\subsection{The $0$-th degeneracy locus and the height inequality}
Let $g\ge 1$. In this and the next section we consider $\mathfrak{A}_g$ as a smooth,
irreducible, quasi-projective variety defined over $\IQbar$.

The $0$-th degeneracy locus  is closely related
to the following height inequality proved by both authors of the
current paper and Dimitrov.


Below we fix a height function $h\colon S(\IQbar)\rightarrow [0,\infty)$
and a fiberwise canonical height $\hat h\colon
\mathfrak{A}_g(\IQbar)\rightarrow[0,\infty)$
as in \cite{DGHUnifML}. 

\begin{theorem}[{\!\!\cite[Theorem~1.6 and Theorem~B.1]{DGHUnifML}}]\label{ThmHtInequality}
  Let $X$ be an irreducible subvariety of $\mathfrak{A}_g$ defined over $\IQbar$. Assume $X_\IC^{\mathrm{deg}}(0)$ is not Zariski dense in $X$. Then there exist a constant $c>0$ and a Zariski open dense subset $U$ of $X$ such that
  \begin{equation}\label{EqHtInequality}
    h(\pi(x)) \le c ( \hat{h}(x) + 1) \quad \text{for all }x \in U(\IQbar).
  \end{equation} 
\end{theorem}

Indeed, the assumption of \cite[Theorem~B.1]{DGHUnifML} is that $X$ is
\textit{non-degenerate} as defined by \cite[Definition~B.4]{DGHUnifML}.
If $X$ fails to be non-degenerate, then 
the mixed Ax--Schanuel Theorem for $\mathfrak{A}_g$
by the first-named author~\cite{GaoMixedAS} implies that
$X_\IC^{\mathrm{deg}}(0)$
contains a non-empty  open subset of $X^{\mathrm{an}}$, 
see \cite[Theorem~1.7]{GaoBettiRank} or 
\cite[Theorem~4.3.1]{GaoHDR}.


\subsection{The $1$-st degeneracy locus and Relative Manin--Mumford}
The $1$-st degeneracy locus %$X^{\mathrm{deg}}(1)$
is closely related to the Relative Manin--Mumford Conjecture. In
 $\mathsection$\ref{SectionConclusion} we will deduce Theorem~\ref{MainThm}
in the case $L=\IQbar$ from the following theorem.

% as an application of the criterion of $X^{\mathrm{deg}}(1) = X$, the relative 
% Manin--Mumford conjecture follows from the following result; see  for the deduction.
%In fact, by
%\cite[Prop.1.10]{GaoBettiRank}, to prove the relative
%Manin--Mumford conjecture (Theorem~\ref{MainThm}) it suffices to prove
%the following result.%, which was stated as \cite[Conj.1.9]{GaoBettiRank}.
\begin{theorem}\label{TheoremFirstDegLocusIsX}
Let $X$ be an irreducible subvariety of $\mathfrak{A}_g$ defined over
$\IQbar$ with $\dim X\ge 1$. If $X(\IQbar) \cap (\mathfrak{A}_g)_{\mathrm{tor}}$ is Zariski dense in $X$, then $X_\IC^{\mathrm{deg}}(1)$ is Zariski dense in $X_\IC$.
\end{theorem}
A large part of this paper
($\mathsection$\ref{SectionLGO}--\ref{SectionAppAS}) is devoted to the
proof of Theorem~\ref{TheoremFirstDegLocusIsX}.



Our proof of Theorem~\ref{TheoremFirstDegLocusIsX} will be arranged as
follows. We divide into two cases: either $X_\IC^{\mathrm{deg}}(0)$ is
Zariski dense in $X$, or $X_\IC^{\mathrm{deg}}(0)$ is not Zariski
dense in $X$. In the first case, the conclusion of
Theorem~\ref{TheoremFirstDegLocusIsX} holds true because
$X_\IC^{\mathrm{deg}}(0) \subseteq X_\IC^{\mathrm{deg}}(1)$. In the second
case, we follow the Pila--Zannier method. More precisely, we will
invoke the height inequality \eqref{EqHtInequality} and a result of 
%David \cite[Th\'eor\`eme~1.4]{DavidMinHaut} 
R\'{e}mond \cite[Proposition~2.9]{RemondGaloisBound} to show that the Galois orbit
of each point in $X(\IQbar) \cap (\mathfrak{A}_g)_{\mathrm{tor}}$ is large,
and then apply the semi-rational variant of the Pila--Wilkie counting
theorem \cite[Corollary~7.2]{HabeggerPilaENS}.
One may also use an older result of David
\cite{DavidMinHaut} to deduce largeness of the Galois orbit, as
explained in the appendix.
We will thus
produce a curve whose projection to the Betti fiber is a
semi-algebraic curve. Finally, by the mixed Ax--Schanuel theorem for
$\mathfrak{A}_g$ (Theorem~\ref{ThmAS}) such a curve will be seen to contribute to
$X_\IC^{\mathrm{deg}}(1)$.



%At this stage, it is worth making some extra comments on Theorem~\ref{ThmDegLocusZarClosed} and \cite[Proposition~1.10]{GaoBettiRank} due to the importance of these two results to the
%current paper. These two results were proved together by a standard
%argument on unlikely intersection problems on $\mathfrak{A}_g$: First
%one shows that the maximal members in the union from
%\eqref{EqDefnDegLocus} are all \textit{weakly optimal subvarieties} as
%introduced by the second-named author and Pila \cite{HabeggerPilaENS},
%then one applies a finiteness result \textit{\`{a} la Bogomolov}
%proved by the first-named author \cite[Theorem 1.4]{GaoMixedAS} to show that the union from \eqref{EqDefnDegLocus} is a finite union.


%%% Local Variables:
%%% TeX-master: "main"
%%% End:



%% Section 4 (new, was 3)
\section{Large Galois Orbits}
\label{SectionLGO}
%!!!TODO: REMOVE PP HYPOTHESIS IN PROP 3.3.

The goal of this section is to prove the following result, claiming
that the Galois orbits of torsion points are \textit{large}.

Let $g\ge 1$. We consider $\mathfrak{A}_g$ as a smooth, irreducible,
quasi-projective variety defined over $\IQbar$.
Let $X$ be an irreducible subvariety of $\mathfrak{A}_g$ defined over
$\IQbar$. Suppose $K\subset\IQbar$ is a number field over which $X$
and $\mathfrak{A}_g$ are defined. For each $x \in X \cap
(\mathfrak{A}_g)_{\mathrm{tor}}$, let $\ord{x}$ denote the order of
$x$.
% and let $N(x)$ be the smallest positive
% integer such that $[N(x)]x$ lies in the zero section of
% $\mathfrak{A}_g \rightarrow \mathbb{A}_g$.
% let $K(x)$ be the smallest
% subfield of $\IQbar$ contained $K$ over which $x$ is defined.

\begin{proposition}\label{PropLGO}
  Assume that $X_\IC^{\mathrm{deg}}(0)$ is not Zariski dense in $X$. Then there exist constants $c>0$, $\delta>0$, and a Zariski open dense subset $U$ of $X$ such that
  \begin{equation}
    \#\mathrm{Gal}(\IQbar/K) x \ge c \cdot \ord{x}^{\delta} \quad \text{ for all }x \in U(\IQbar) \cap (\mathfrak{A}_g)_{\mathrm{tor}}.
  \end{equation}
\end{proposition}

The key ingredients of the proof are the height inequality
\cite[Theorems 1.6 and B.1]{DGHUnifML} and a \textit{quantitative} version of Masser's work on  the Galois orbit of torsion points on abelian varieties. In the proof presented here, we use R\'{e}mond's \cite[Proposition~2.9]{RemondGaloisBound}, which
is based on
%Masser and W\"{u}stholz's deep Isogeny Theorem \cite{MW:abelianisog}
%and its refinement by
earlier work of 
Gaudron--R\'{e}mond \cite{GR:Isogeny} itself related to isogeny
estimates of  Masser--W\"{u}stholz~\cite{MW:abelianisog}.
Alternatively, one can use \cite[Corollaire~1.7]{GRGaloisBound}.
In the Appendix, we show how an older result of David \cite{DavidMinHaut} suffices for our purpose. %David's result does not rely on Masser--W\"{u}stholz.


%The proof of Proposition~\ref{PropDavidBound} 

%\subsection{Proof of Proposition~\ref{PropLGO}}

\begin{proof}
Let $U$ be the Zariski open dense subset of $X$ obtained from
Theorem~\ref{ThmHtInequality}, and let $c = c(X) >0$ be the constant
from the same theorem. Then
\begin{equation}\label{Eq}
  h(\pi(x)) \le c \qquad \text{for all }x \in U(\IQbar) \cap (\mathfrak{A}_g)_{\mathrm{tor}}
\end{equation}
as the N\'eron--Tate height of a torsion point is $0$.

% The subset $S_0= \{s \in S(\IQ) : h(s) < c\}$, is a finite set by the Northcott property. Thus $U \setminus \pi^{-1}(S_0)$ is again Zariski open dense in $X$. Replacing $U$ by $U \setminus \pi^{-1}(S_0)$, we may and do assume that $\IQ(x) \not= \IQ$ for all $x \in U(\IQbar) \cap (\mathfrak{A}_g)_{\mathrm{tor}}$. In particular, $[\IQ(x):\IQ] \ge 2$ for all $x \in U(\IQbar) \cap (\mathfrak{A}_g)_{\mathrm{tor}}$.

For $x\in U(\IQbar)\cap (\mathfrak{A}_g)_{\mathrm{tor}}$
the abelian variety $A=\mathfrak{A}_{g,x}$
is  defined over $k=K(x)$ with $K$ as in the beginning of
this section. %$\mathsection$\ref{SectionLGO}. 

Let $h_{\mathrm{Fal}}(A)$ denote the \textit{stable Faltings height} of
$A$. The additive normalization in $h_{\mathrm{Fal}}$ plays no role in
the current work. 
Because of \eqref{Eq} and
by the fundamental work of Faltings~\cite[$\mathsection 3$ including
 the proof of Lemma~3]{Faltings:ES} (see also
\cite[the remarks below Proposition~V.4.4 and Proposition~V.4.5]{FaltingsChai}), we have $h_{\mathrm{Fal}}(A) \le c'$ for some constant $c' = c'(X)>0$.% Up to enlarging $c'$, we may and do assume $c' > 1$.

%Next we invoke a quantitative version of Masser's work on the Galois orbit of torsion points on abelian varieties
Define
\[
\kappa(X) = ((14g)^{64g^2} [k:\IQ] \max\{1,c', \log [k:\IQ]\}^2)^{1024 g^3}.
\]
Then by R\'{e}mond   \cite[Proposition~2.9]{RemondGaloisBound}, we have $\ord{x} \le \kappa(X)^{4g+1}$ because $h_{\mathrm{Fal}}(A) \le c'$. Hence we are done.
\end{proof}

%%% Local Variables:
%%% TeX-master: "main"
%%% End:








%% Section 5
\section{Point Counting}
Let $g\ge 1$ be an integer. In this and the next section we consider
$\mathfrak{A}_g$ as a smooth, irreducible, quasi-projective variety
defined over $\IQbar$.
Let $X$ be an irreducible subvariety of $\mathfrak{A}_g$, as usual defined over
$\IQbar$. Let $K$ be a number field such that both $\mathfrak{A}_g$ and $X$
are defined over $K$. Recall the uniformization $u \colon \IC^g \times \mathfrak{H}_g
\rightarrow \mathfrak{A}_g^{\mathrm{an}}$ from
\eqref{EqUnifUniversalAbVar}.
Let $\tilde{b} \colon \IC^g \times \mathfrak{H}_g \rightarrow \IR^{2g}$ be the universal uniformized Betti map from \eqref{EqUnivBettiMap}; it is a semi-algebraic map.%Let $p \colon \IC^g \times \mathfrak{H}_g \rightarrow \IC^g$ be the natural projection.

The goal of this section is to prove the following result. Let $\IRanexp$ denote the  structure generated by the unrestricted
exponential function on the real numbers and the restriction of all
real analytic functions to a hypercube. Then $\IRanexp$ is o-minimal by the
Theorem of van den Dries and Miller \cite{DriesMiller:94} and earlier work of Wilkie \cite{Wilkie:96}.
Throughout this section, the word \textit{definable} will mean definable in $\IRanexp$. 


\begin{proposition}\label{PropositionSemiAlgCounting}
  Assume $\dim X\ge 1$ and
  that $X_\IC^{\mathrm{deg}}(0)$ is not Zariski dense in $X$. Then there exists a Zariski open
  dense subset $U$ of $X$ and $B\ge 1$ with the following property. For each $x \in
  U(\IQbar) \cap (\mathfrak{A}_g)_{\mathrm{tor}}$ with
  $\mathrm{ord}(x)\ge B$
  there exist $\tilde\gamma_x \colon
  [0,1]\rightarrow\IC^g\times \mathfrak{H}_g$ continuous and definable with
  the following properties. % Let $\gamma_x = u\circ\tilde\gamma_x
  % \colon [0,1]\rightarrow \mathfrak{A}_g$, then 
  % a continuous non-constant map $\gamma_x \colon [0,1] \rightarrow X^{\mathrm{an}}$ such that
  \begin{enumerate}
  \item[(i)] The map $\tilde\gamma_x$ is non-constant,
    $\tilde\gamma_x|_{(0,1)}$ is real analytic,  and $u(\tilde\gamma_x([0,1]))\subset
    X^{\mathrm{an}}$, 
  \item[(ii)] $u(\tilde\gamma_x(0)) \in \mathrm{Gal}(\IQbar/K)x$, and
  \item[(iii)] $\tilde{b} \circ \tilde{\gamma}_x\colon
    [0,1]\rightarrow\IR^{2g}$ is semi-algebraic. % , where  $\tilde{\gamma}_x \colon [0,1] \rightarrow \IC^g \times \mathfrak{H}_g$ is a continuous lift of $\gamma_x$ for the uniformization $u \colon \IC^g \times \mathfrak{H}_g \rightarrow \mathfrak{A}_g^{\mathrm{an}}$.
    % \item[(iii)] $\pi \circ \gamma$ is non-constant.
  \end{enumerate}
\end{proposition}




\begin{proof} Let $U$ be the Zariski open dense subset of $X$ from Proposition~\ref{PropLGO}.

  %% Removed the next two paragraphs.

  % Observe that $X$ is not entirely contained in  $\ker [N]$
  % for some integer $N\ge 1$.
  % Indeed, otherwise and as $\dim X\ge 1$ all of $X$ itself would
  % appear as some $Y$ in the union (\ref{EqDefnDegLocus}) for $t=0$.
  % But this would entail $X_\IC^{\mathrm{deg}}(0)=X$, contrary to the
  % hypothesis.

  % Thus we may remove  subsets $\ker [N]|_X$, for finitely many $N$, from $U$
  % and
  % retain a Zariski open and dense subset of $X$. 
  
Let $\tau \colon \IR^{2g} \times \mathfrak{H}_g \rightarrow \IC^g \times \mathfrak{H}_g$ be the semi-algebraic isomorphism given by \eqref{EqComplexStrOfX2g}. We have the following commutative diagram.%  This map is important in the proof. . These two maps and the natural projection $p$ introduced above Proposition~\ref{PropositionSemiAlgCounting} fit into the following commutative diagram:
\[
\xymatrix{
\IC^g \times \mathfrak{H}_g \ar[rd]_-{\tilde{b}} & \IR^{2g} \times \mathfrak{H}_g \ar[l]_-{\tau} \ar[d] \\
 & \IR^{2g}
}
\]






By work of Peterzil--Starchenko
\cite{PeterzilStarchenkoDefinabilityTheta}, there exists a
semi-algebraic fundamental subset $\mathfrak{F}_0 \subseteq
\mathfrak{H}_g$ with the following property: for each integer $M > 0$,
if we set $\mathfrak{F} = \tau([-M,M]^{2g} \times \mathfrak{F}_0)
\subseteq \tau(\IR^{2g} \times \mathfrak{H}_g) = \IC^g \times
\mathfrak{H}_g$, then $u|_{\mathfrak{F}} \colon \mathfrak{F}
\rightarrow \mathfrak{A}_g^{\mathrm{an}}$ is  definable in the
o-minimal structure $\IRanexp$;
observe that we may assume that $\mathfrak{A}_g$ is embedded in some
projective space. %As an upshot, $(u |_{\mathfrak{F}})^{-1}(X)$ is a definable subset of $\IC^{g} \times \mathfrak{H}_g$. 
Note that $\mathfrak{F}$ is a semi-algebraic subset of $\IC^g \times \mathfrak{H}_g$ because both $\mathfrak{F}_0$ and $\tau$ are semi-algebraic.

Fix an integer $M \ge 1$ large enough such that $u(\mathfrak{F}) = \mathfrak{A}_g^{\mathrm{an}}$.

%Observe that for any  $x \in X \cap (\mathfrak{A}_g)_{\mathrm{tor}}$,  we have $\tilde{b}(\tilde{x}) \in \IQ^{2g}$ for each $\tilde{x} \in u^{-1}(x)$. Moreover if $N(x) \in \IZ$ such that $[N(x)]x$ is in the zero section, then $N(x) \tilde{b}(\tilde{x}) \in \IZ^{2g}$.



Consider $\IR^{2g} \times \mathfrak{F} \subseteq \IR^{2g} \times \IC^g \times \mathfrak{H}_g$. Define
\begin{equation}\label{EqDefnDefFamily}
\hat{X} = \{(r, \tilde{x}) \in \IR^{2g} \times \mathfrak{F} : \tilde{x} \in (u |_{\mathfrak{F}})^{-1}(X(\IC)), ~ r = \tilde{b}(\tilde{x})\}.
\end{equation}
%Then $\hat X$ is definable in $\IRanexp$. 
As both $u|_{\mathfrak{F}}$ and $\tilde{b}$ are definable maps, the set $\hat{X}$ is a definable subset of $\IR^{2g} \times \mathfrak{F}$. By choice of $\mathfrak{F}$, the following holds true:  for any $(r,\tilde{x}) \in \hat{X}$, we have $r \in [-M,M]^{2g}$.

The height $H(p/q)$ of a rational number $p/q$ with coprime integers
$p,q$ and $q\ge 1$ is $\max\{|p|,q\}$. The height $H(r_1,\ldots,r_n)$
with $r_1,\ldots,r_n\in \IQ$ is $\max_i H(r_i)$. 

For each $T \ge 1$, set $\hat{X}(T) = \{(r,\tilde{x}) \in \hat{X} : r \in \IQ^{2g},~ H(r) \le T\}$. 

Now let $x \in X(\IQbar) \cap (\mathfrak{A}_g)_{\mathrm{tor}}$ and let
$T=\ord{x}$ be the order of $x$. 


Choose $\tilde{x} \in (u|_{\mathfrak{F}})^{-1}(x)$ and set $r = \tilde{b}(\tilde{x})$. Then $r \in \IQ^{2g}$ and $T\cdot r \in \IZ^{2g}$. As $r \in [-M,M]^{2g}$ by choice of $\mathfrak{F}$, we have $H(r) \le MT$. So $(r,\tilde{x}) \in \hat{X}(MT)$.


Thus we can invoke Proposition~\ref{PropLGO} to obtain a subset $\Sigma \subseteq \hat{X}(MT)$ as follows. Recall that $X$ is defined over a number field $K$, and consider $\mathrm{Gal}(\IQbar/K)x \subseteq X(\IQbar)$. Let $c$ and $\delta$ be the constants from Proposition~\ref{PropLGO}, then we have $\#\mathrm{Gal}(\IQbar/K)x \ge c T^{\delta}$. % up to replacing $c$ by $c/[K:\IQ]$.
Each point $x' \in \mathrm{Gal}(\IQbar/K)x$ is again in $X(\IQbar)
\cap (\mathfrak{A}_g)_{\mathrm{tor}}$ of order $T$. Hence it gives rise to a point $(r',\tilde{x}') \in \hat{X}(MT)$. Let $\Sigma$ be the set of all such points in $\hat{X}(MT)$. Then for the natural projection $\pi_2 \colon \IR^{2g}\times \mathfrak{F} \rightarrow \mathfrak{F}$, we have
\begin{equation}
 \#\pi_2(\Sigma) \ge cT^{\delta}. % > c'(MT)^{\delta/2}
\end{equation}
We may fix $B$ large in terms of $c,M,$ and $c'$, the constant in
the semi-rational variant of the Pila--Wilkie
counting theorem \cite[Cor.7.2]{HabeggerPilaENS} with $\epsilon =
\delta/2$, to the following effect.
If  $T = \ord{x}\ge B$, then
\begin{equation}
 \#\pi_2(\Sigma)  > c'(MT)^{\delta/2}.
\end{equation}

% \philipp{Careful, only $\#\Sigma$ is large, $\#\pi_2(\Sigma)$ could be
%   small, say $<T^{\delta/2}$. But then, by the pigeonhole principle,
%   there would be some $Z\in\mathfrak{F}$.}

%%% The following paragraph doesn't fit here. We can take $B$ large.
% As described near the beginning of this proof, we may assume that $U$
% contains no torsion points of order from a finite prescribed set.
% Therefore, we may assume that $T$ is large enough as described above.

We apply \cite[Cor.7.2]{HabeggerPilaENS} to $\hat X$, taken as a
single-fiber family, $\Sigma \subseteq \hat{X}(MT)$, and $\epsilon =
\delta/2$. We obtain a continuous and definable map $\beta \colon
[0,1] \rightarrow \hat{X}$ such that
\begin{enumerate}
\item[(a)] the composition $[0,1]\xrightarrow{\beta} \hat{X} \subseteq \IR^{2g} \times \mathfrak{F} \xrightarrow{\pi_1} \IR^{2g}$ is semi-algebraic;
\item[(b)] the composition $\pi_2 \circ \beta \colon [0,1]\rightarrow \hat{X} \rightarrow \mathfrak{F}$ is non-constant;
\item[(c)] $\pi_2(\beta(0)) \in \pi_2(\Sigma)$;
\item[(d)] $\beta|_{(0,1)}$ is real analytic. 
\end{enumerate}
For the final property we require that $\IRanexp$ has analytic
cell decomposition, this follows from the
work of van den Dries and Miller \cite{DriesMiller:94}. %vdDries-Miller 5.1(3)

Let $\tilde\gamma_x = \pi_2\circ\beta\colon [0,1]\rightarrow
\IC^g\times \mathfrak{H}_g$. 
We will show that this $\tilde\gamma_x$ is what
we desire.


% Let $\gamma_x = u\circ \pi_2\circ\beta \colon [0,1] \rightarrow \mathfrak{A}_g^{\mathrm{an}}$. We will show that this $\gamma_x$ is what we desire.

First, $u(\tilde\gamma_x([0,1])) = u(\pi_2\circ \beta([0,1])) \subseteq
u(\pi_2(\hat{X}))\subset X(\IC)$. 
% \subseteq u((u |_{\mathfrak{F}})^{-1}(X)) \subseteq
% X$
 Moreover, $\tilde\gamma_x$ is non-constant by property (b) above
and its restriction to $(0,1)$ is real analytic by (d). So part (i) follows.

By property (c), we have $u(\tilde\gamma_x(0)) = (u\circ \pi_2\circ\beta)(0) \in u(\pi_2(\Sigma)) \subseteq \mathrm{Gal}(\IQbar/K)x$. Hence (ii) is established.

It remains to establish property (iii). Notice that
$\pi_1|_{\hat X} = \tilde b \circ \pi_2|_{\hat X}$ by
(\ref{EqDefnDefFamily}).
As $\beta$ takes values in $\hat X$ we find
$\pi_1\circ\beta = \tilde b \circ\pi_2\circ\beta = \tilde b
\circ\tilde \gamma_x$. So $\tilde b \circ\tilde\gamma_x$ is
semi-algebraic by (a). 
% Notice that $\pi_2\circ \beta$ is by construction a lift of $\gamma_x$. Thus we can set $\tilde{\gamma}_x = \pi_2 \circ  \beta$. Therefore $\tilde{b} \circ \tilde{\gamma}_x = \tilde{b} \circ \pi_2 \circ \beta$. On the other hand, we have $\pi_1|_{\hat{X}} = \tilde{b} \circ \pi_2|_{\hat{X}}$ by definition of $\hat{X}$ \eqref{EqDefnDefFamily}. Thus $\tilde{b} \circ \tilde{\gamma}_x = \pi_1 \circ \beta$. So $\tilde{b} \circ \tilde{\gamma}_x$ is semi-algebraic by property (a) above, and we are done. 
\end{proof}

%%% Local Variables:
%%% TeX-master: "main"
%%% End:


%% Section 6
\section{Application of Mixed Ax--Schanuel}
\label{SectionAppAS}
Let $g\ge 1$. In this and the next section we consider
$\mathfrak{A}_g$ as a smooth, irreducible, quasi-projective variety
defined over $\IQbar$.

\begin{proposition}\label{PropXIsDeg1}
  Let $X$ be an irreducible subvariety of
  $\mathfrak{A}_g$ and assume that
  $\dim X\ge 1$, that $X_\IC^{\mathrm{deg}}(0)$ is not Zariski dense
  in $X$, and that $X(\IQbar)\cap
  (\mathfrak{A}_g)_{\mathrm{tor}}$ is Zariski dense in $X$.
  Then $X_\IC^{\mathrm{deg}}(1)$ is Zariski dense in $X_\IC$.
\end{proposition}

\begin{proof}
  Let $K$ be a number field over which $\mathfrak{A}_g$ and $X$ are
  defined.


  Let $U\subset X$ and $B\ge 1$ be 
  from Proposition~\ref{PropositionSemiAlgCounting}. We may apply
  Proposition~\ref{PropositionSemiAlgCounting} to elements of $\Sigma=
  \{ x\in U(\IQbar) \cap (\mathfrak{A}_g)_{\mathrm{tors}} :
  \mathrm{ord}(x) \ge B\}$. By replacing $U$ by a Zariski open and dense
  subset we may assume that $\Sigma$ is stable under the action of
  $\mathrm{Gal}(\IQbar/K)$.

  Suppose that $\Sigma$ is not Zariski dense in $X$. Then by hypothesis $X$ contains a Zariski
  dense set of torsion points of order  $<B$. Then $X \subset
  \mathrm{ker}[n]$ for some positive integer $n< B$. As $\dim X >0$ we
  find that $X_\IC$ appears in the union (\ref{EqDefnDegLocus}) for $t=0$. In particular,  
  $X_\IC^{\mathrm{deg}}(0)=X$ which contradicts our hypothesis.
  So $\Sigma$ is Zariski dense in $X$. 
  % Because $X^{\mathrm{deg}}(1)$ is Zariski closed in $X$ and by the assumption $(X\cap (\mathfrak{A}_g)_{\mathrm{tor}})^{\mathrm{Zar}} = X$, to prove Proposition~\ref{PropXIsDeg1} it suffices to show that
  % \begin{equation}\label{EqXIsDeg1}
  %   U \cap (\mathfrak{A}_g)_{\mathrm{tor}} \subseteq X^{\mathrm{deg}}(1).
  % \end{equation}

% By Theorem~\ref{ThmDegLocusZarClosed}, $X^{\mathrm{deg}}(1)$ is
% Zariski closed in $X$. 
% By Theorem~\ref{ThmDegLocusZarClosed}, $X^{\mathrm{deg}}(1)$ is Zariski closed and is defined over $\IQbar$. 
% Let $K$ be a number field such that both $X$ and $X^{\mathrm{deg}}(1)$
% are defined over $K$.
% \philipp{Do we really need to know that $X^{\mathrm{deg}}(1)$ is over
%   $K$? It should follow from the arguments below.}

  Take $x \in \Sigma$ and % $U(\IQbar) \cap (\mathfrak{A}_g)_{\mathrm{tor}}$.
  let $\tilde\gamma_x$ be as in Proposition~\ref{PropositionSemiAlgCounting}.
%   For the
% uniformization $u \colon \IC^g \times \mathfrak{H}_g \rightarrow
% \mathfrak{A}_g^{\mathrm{an}}$, the non-constant continuous map
% $\gamma_x \colon [0,1] \rightarrow X^{\mathrm{an}}$ from
% Proposition~\ref{PropositionSemiAlgCounting}
% equals $u\circ\tilde\gamma_x$ with 
% $\tilde{\gamma}_x \colon [0,1] \rightarrow \IC^g\times\mathfrak{H}_g$.

  Let $\tilde{C} = \tilde\gamma_x([0,1]) \subset\IC^g\times\mathfrak{H}_g$. This is a connected and
  definable set as $\tilde\gamma_x$ is continuous and definable;
  recall that definable means definable in $\IRanexp$.
  Moreover, $\dim\tilde{C} \ge 1$ since $\tilde\gamma_x$ is non-constant.

  We let $\tilde Z$ denote the intersection of all complex analytic subsets
  of $\IC^g\times \mathfrak{H}_g$ containing $\tilde C$. Then $\tilde Z$ is
  itself a complex analytic subset of $\IC^g\times\mathfrak{H}_g$.
  Moreover, $\tilde Z$ is irreducible as $\tilde\gamma_x|_{(0,1)}$ is real
  analytic. Finally, $\dim \tilde Z\ge 1$.

  We define $Y\subset\mathfrak{A}_{g,\IC}$
  to be the Zariski closure $u(\tilde Z)^{\mathrm{Zar}}$. Then $Y$ is
  irreducible and $\dim Y\ge 1$. 
  % Moreover, $u(\tilde C) \subset \mathfrak{A}_{g}(\IC)$
  % is connected and infinite. 
  % We define $Y$ to be  $u(\tilde{C})^{\mathrm{Zar}}$, taking the
  % Zariski closure in $\mathfrak{A}_{g,\IC}$. 
  % Then $Y$ is irreducible since $u\circ\tilde \gamma_x|_{(0,1)}$ is
  % real analytic. Moreover, $\dim Y\ge 1$. 


  Next we apply the first-named author's mixed Ax--Schanuel Theorem for the
  universal family to obtain the following lemma. 
  \begin{lemma}
    \label{LemmaZarClosureinXdeg1}
    We have $Y \subseteq X_\IC^{\mathrm{deg}}(1)$.
  \end{lemma}

  Let us finish the proof of Proposition~\ref{PropXIsDeg1} assuming
  Lemma~\ref{LemmaZarClosureinXdeg1}. By property (ii) of
  Proposition~\ref{PropositionSemiAlgCounting}, we have $x' \in
  u(\tilde{C})$  for some $x' \in \mathrm{Gal}(\IQbar/K)x$. Thus $x' \in
  Y(\IC)$. Assuming Lemma~\ref{LemmaZarClosureinXdeg1}, we then have $x'
  \in X_\IC^{\mathrm{deg}}(1)$.

  To summarize, the orbit under $\mathrm{Gal}(\IQbar/K)$ of any given
  element in $\Sigma$ meets $X_\IC^{\mathrm{deg}}(1)$. Recall that
  $\mathrm{Gal}(\IQbar/K)$ acts on 
  $\Sigma$.
  So there is a subset $\Sigma'\subset \Sigma$ contained
  completely in $X_\IC^{\mathrm{deg}}(1)$ such that for all $x\in \Sigma$
  there is $\sigma\in \mathrm{Gal}(\IQbar/K)$ with $\sigma(x)\in
  \Sigma'$.
  The Zariski closure of $\Sigma'$ in $X_\IC$ is defined over $\IQbar$.
  So it is stable under a finite index subgroup of
  $\mathrm{Gal}(\IQbar/K)$.
  As $\Sigma$ is Zariski dense in $X$ we conclude that $\Sigma'$ is
  Zariski dense in $X$ as well.

  As $\Sigma' \subset X_\IC^{\mathrm{deg}}(1)$ we conclude that
  $X_\IC^{\mathrm{deg}}(1)$ is Zariski dense in $X_\IC$, as desired. 
  % The argument above applies to each $x \in U(\IQbar) \cap
% (\mathfrak{A}_g)_{\mathrm{tor}}$. So $U(\IQbar) \cap
% (\mathfrak{A}_g)_{\mathrm{tor}} \subseteq
% X^{\mathrm{deg}}(1)(\IQbar)$. But $U$ is Zariski open and dense in $X$
% and $(X\cap (\mathfrak{A}_g)_{\mathrm{tor}})^{\mathrm{Zar}} = X$. So
% taking Zariski closures we obtain that $X$ is the Zariski closure of $X^{\mathrm{deg}}(1)$. Hence we are done.
\end{proof}

\begin{proof}[Proof of Lemma~\ref{LemmaZarClosureinXdeg1}]
  We systematically work with varieties defined over $\IC$ and  also identify
  varieties with their set of complex points.
  To ease notation we write $X$ for the base change $X_\IC$.


  We apply the weak Ax--Schanuel theorem for $\mathfrak{A}_g$,
  Theorem~\ref{ThmAS}, to $\tilde Z$ and get
  \begin{equation*}
    \dim {\tilde Z}^{\mathrm{Zar}} + \dim u(\tilde Z)^{\mathrm{Zar}} \ge \dim \tilde Z + \dim u(\tilde Z)^{\mathrm{biZar}}.
  \end{equation*}
  Recall $\dim \tilde Z\ge 1$ and  $Y=u(\tilde
  Z)^{\mathrm{Zar}}\subset u(\tilde Z)^{\mathrm{biZar}}$. So
  $Y^{\mathrm{biZar}}\subset  u(\tilde Z)^{\mathrm{biZar}}$
  and hence
  \begin{equation}
    \label{eq:masineq}
    \dim {\tilde Z}^{\mathrm{Zar}} + \dim Y \ge 1 + \dim Y^{\mathrm{biZar}}.
  \end{equation}  

  Consider the following commutative (possibly non-Cartesian) diagram involving uniformizations and projections to the bases:
  \begin{equation*}%\label{EqDiagUnifProj}
    \xymatrix{
      \IR^{2g} & \IC^g \times \mathfrak{H}_g \ar[l]_-{\tilde{b}} \ar[r]^-{u} \ar[d]_{\tilde{\pi}} & \mathfrak{A}_g^{\mathrm{an}} \ar[d]^{\pi} \\
      & \mathfrak{H}_g \ar[r]^-{u_B} & \mathbb{A}_g^{\mathrm{an}}.
    }
  \end{equation*}
  Recall that  $\tilde b$
  % \[
  %   \tilde{b} \colon \IC^g \times \mathfrak{H}_g \rightarrow \IR^{2g}
  % \]
  is the universal uniformized Betti map defined in
  \eqref{EqUnivBettiMap}.

  We have
  \begin{equation*}
    u_B(\tilde \pi(\tilde C)) = \pi(u(\tilde C)) \subset
    \pi(u(\tilde Z)) \subset \pi(u(\tilde Z)^{\mathrm{Zar}})=\pi(Y) \subset \pi(Y)^{\mathrm{biZar}}.
  \end{equation*}
  So $\tilde\pi\circ\tilde \gamma_x$ takes values in
  $u_B^{-1}(\pi(Y)^{\mathrm{biZar}})$. As
  $\tilde\pi\circ\tilde\gamma_x|_{(0,1)}$ is real analytic and
  continuous at the boundary, the values
  of $\tilde\pi\circ\tilde \gamma_x$ lie in an complex analytic
  irreducible component $\tilde Y_0$ of   $u_B^{-1}(\pi(Y)^{\mathrm{biZar}})$. By
  the definition of bi-algebraic subsets of $\IA_g$ there is an
  algebraic subset $\tilde Y$ of $\mathrm{Mat}_{g\times g}(\IC)$ such that
  $\tilde Y_0$ is a complex analytic irreducible component of
  $\tilde Y \cap \mathfrak{H}_g$.  We may assume that $\tilde Y$ is
  irreducible.
  In particular,
  \begin{equation}
    \label{eq:dimYequality}
   \dim \tilde Y = \dim \tilde Y_0 =
  \dim \pi(Y)^{\mathrm{biZar}}. 
  \end{equation}

  Recall that $\tilde\pi(\tilde C)=\tilde\pi(\tilde\gamma_x([0,1]))
  \subset \tilde Y_0$. So
  $$
  \tilde C \subset\IC^g \times \tilde Y_0 \subset\IC^g\times \tilde Y.
  $$

  On the other hand, $\tilde b\circ\tilde\gamma_x$ is semi-algebraic
  by Proposition~\ref{PropositionSemiAlgCounting}(iii). In other
  words, the Betti coordinates along $\tilde\gamma_x$ lie in a
  real semi-algebraic curve in $\IR^{2g}$. Thus there is a complex algebraic curve
  $A\subset\IC^{2g}$ with $\tilde C\subset \{((\tau,1_g)z,\tau) : z\in
  A\text{ and }\tau\in\mathfrak{H}_g\}$ (we take points of $A$ as
  column vectors).
  We conclude $\tilde C\subset E$ with $E$ the Zariski closure of 
  \begin{equation*}
    (\IC^g\times \tilde Y)\cap \{((\tau,1_g)z,\tau) : z\in
    A\text{ and }\tau\in\mathrm{Mat}_{g\times g}(\IC)\}.
  \end{equation*}
  
  Recall that $\tilde Z$ is the smallest complex analytic subset of
  $\IC^g\times \mathfrak{H}_g$ containing
  $\tilde C$. So it is contained in $E$. Moreover,  $E$ is complex
  algebraic, hence
    ${\tilde Z}^{\mathrm{Zar}} \subset E$.
  Finally, all fibers of the restricted projection $E\rightarrow
  \mathrm{Mat}_{g\times g}(\IC)$ have dimension at most $1$.
  Therefore, $\dim E\le 1+\dim \tilde Y$ and hence
  $\dim {\tilde Z}^{\mathrm{Zar}}\le \dim E \le \dim 1 + \dim \tilde Y = 1 +\dim
  \pi(Y)^{\mathrm{biZar}}$ by (\ref{eq:dimYequality}). 

  We apply (\ref{eq:masineq}) and conclude
  $\dim Y \ge \dim Y^{\mathrm{biZar}} - \dim \pi(Y)^{\mathrm{biZar}}$.
  As $\dim Y\ge 1$ we find $Y\subset X^{\mathrm{deg}}(1)$ by
  (\ref{EqDefnDegLocus}).
  This concludes the proof of Lemma~\ref{LemmaZarClosureinXdeg1}.
\end{proof}

\subsection{Proof of Theorem~\ref{TheoremFirstDegLocusIsX}}
Let $X$ be an irreducible subvariety of $\mathfrak{A}_g$ defined over $\IQbar$ and $\dim X \ge 1$. Assume that $X(\IQbar) \cap (\mathfrak{A}_g)_{\mathrm{tor}}$ is Zariski dense in $X$.

We are in two cases: either $X_\IC^{\mathrm{deg}}(0)$ is or is not Zariski dense
 in $X$, for $X_\IC^{\mathrm{deg}}(0)$ defined by
\eqref{EqDefnDegLocus} (with $t=0$).

If $X_\IC^{\mathrm{deg}}(0)$ is Zariski dense in $X$, so is $X_\IC^{\mathrm{deg}}(1)$ because $X_\IC^{\mathrm{deg}}(1) \supseteq X_\IC^{\mathrm{deg}}(0)$ by definition of $X_\IC^{\mathrm{deg}}(t)$.% But $X^{\mathrm{deg}}(1)$ is Zariski closed in $X$; see Theorem~\ref{ThmDegLocusZarClosed}. So $X = X^{\mathrm{deg}}(1)$.

If $X_\IC^{\mathrm{deg}}(0)$ is not Zariski dense in $X$, then we can
invoke Proposition~\ref{PropXIsDeg1} to conclude that
$X_\IC^{\mathrm{deg}}(1)$ is Zariski dense in $X_\IC$.
%So $X_\IC^{\mathrm{deg}}(1)=X$ by Theorem~\ref{ThmDegLocusZarClosed}.
Hence we are done.\qed

%%% Local Variables:
%%% TeX-master: "main"
%%% End:


%% Section 7
\section{The Relative Manin--Mumford Conjecture over $\IQbar$}
\label{SectionConclusion}
Now we are ready to prove  our main theorem, Theorem~\ref{MainThm}, when the base field is $L=\IQbar$,
using Theorem~\ref{TheoremFirstDegLocusIsX}.
% and the criterion for $X^{\mathrm{deg}}(1) = X$.
The deduction, which uses the criterion for $X_{\IC}^{\mathrm{deg}}(1) = X$, was given by \cite[Proposition 1.10]{GaoBettiRank} and with more details 
by \cite[Corollary 6.3]{GH:nondeg} for the universal family of
principally polarized abelian varieties.
To make the current paper more self-contained, we represent and simplify this proof.%\cite[Prop.1.10]{GaoBettiRank}.% and present more details.
\begin{proof}[Proof of Theorem~\ref{MainThm} over $\IQbar$]
Let us first reduce to the case where $\cA \rightarrow S$ satisfies the following assumption:
\begin{center}
  %\label{def:hyp}
  {\tt (Hyp):} $S$ is a regular locally closed subvariety of $\mathbb{A}_g$ defined over $\IQbar$ and $\cA = \mathfrak{A}_g \times_{\mathbb{A}_g} S$.
\end{center}
Indeed, let $X \subseteq \cA \rightarrow S$ be defined over $\IQbar$ satisfying the assumptions of Theorem~\ref{MainThm}, \textit{i.e.}, %$X$ is dominant to $S$, 
$\mathbb{Z}X$ is Zariski dense in $\cA$ and $X(\IQbar) \cap \cA_{\mathrm{tor}}$ is Zariski dense in $X$. Take a relatively ample line bundle on $\cA \rightarrow S$. By \cite[$\mathsection$2.1]{GenestierNgo}, it induces a polarization of type $D = (d_1,\ldots,d_g)$ on $\cA\rightarrow S$ with $d_1|\cdots|d_g$. Take $\ell \gg 1$ large enough with $(\ell, d_g) = 1$. Then there exists a finite \'{e}tale morphism $\rho \colon S' \rightarrow S$ such that $\cA' := \cA\times_S S' \rightarrow S'$ has level-$\ell$-structure. Write $\rho_{\cA} \colon \cA' \rightarrow \cA$ for the natural projection. Let $X'$ be an irreducible component of $\rho_{\cA}^{-1}(X)$; then $\dim X' = \dim X$. It is not hard to show that 
%$X'$ dominates $S'$, 
$\mathbb{Z}X'$ is Zariski dense in $\cA'$ and $X'(\IQbar) \cap \cA'_{\mathrm{tor}}$ is Zariski dense in $X'$. %  If one can prove Theorem~\ref{MainThm} for $X' \subseteq \cA' \rightarrow S'$, \textit{i.e.}, $\dim X' \ge g$, then we get $\dim X = \dim X' \ge g$. Thus we may assume that $\cA \rightarrow S$ carries a polarization of type $D$ and has level-$\ell$-structure.

Next we have a Cartesian diagram
\[
\xymatrix{
\cA' \ar[r]^-{\iota'} \ar[d] \pullbackcorner & \mathfrak{A}_g \ar[d] \\
S' \ar[r]^-{\iota'_S} & \mathbb{A}_g.
}
\]
Let $S''$ be a regular, Zariski open and dense subset of
$\iota'_S(S')^{\mathrm{Zar}}$ that is contained in the image of $S'$.
Let $\cA'' := \mathfrak{A}_g \times_{\mathbb{A}_g} S''$ and $X'' := \iota'(X') \cap \cA''$. Then $\dim X' \ge \dim X''$. It is not hard to check that 
$\mathbb{Z}X''$ is Zariski dense in $\cA''$, and $X''(\IQbar) \cap
\cA''_{\mathrm{tor}}$ is Zariski dense in $X''$. \textit{If}
Theorem~\ref{MainThm} holds under the extra assumption {\tt (Hyp)}
then  it holds true for $X'' \subseteq \cA'' \rightarrow S''$. So $\dim X'' \ge g$. Hence $\dim X = \dim X' \ge \dim X'' \ge g$, and we are done.

Therefore, from now on without loss of generality we work with $\cA
\rightarrow S$ satisfying {\tt (Hyp)}. % We now prove
% Theorem~\ref{MainThm} for $L=\IQbar$ under this assumption by induction on $g$. 

Let $X \subseteq \cA$ be an irreducible subvariety defined over
$L=\IQbar$ satisfying the assumptions of Theorem~\ref{MainThm} and {\tt (Hyp)}. Our goal is to prove $\dim X\ge g$. We proceed by
induction on $g\ge 1$. 

First observe that  $\dim X \ge 1$. Indeed, if $X$ were a point, it
would be torsion. But this contradicts the hypothesis of
Theorem~\ref{MainThm}.
We thus recover the case $g=1$. 

% \philipp{TODO: Use $X_\IC^{\mathrm{deg}}(t)$ and not
%   $X^{\mathrm{deg}}(t)$.}
We assume $g\ge 2$ and
that Theorem~\ref{MainThm} holds true for all $\cA \rightarrow
S$ satisfying {\tt (Hyp)} of relative dimension
$1,\ldots,g-1$. % We wish to prove it when $\cA \rightarrow S$ has
% relative dimension $g\ge 2$ and satisfies {\tt (Hyp)}.

%<<<<<<< HEAD

By Theorem~\ref{TheoremFirstDegLocusIsX}, $X_\IC^{\mathrm{deg}}(1)$ is
Zariski dense in $X_\IC$. So $X_\IC^{\mathrm{deg}}(1) = X_\IC$ by the Zariski
closedness of $X_\IC^{\mathrm{deg}}(1)$,
Theorem~\ref{ThmDegLocusZarClosed}. Now we apply
\cite[Theorem~8.1]{GaoBettiRank} with $t = 1$. %\philipp{Added the
%  following comment. I think it's important to emphasize where the
%  hypothesis on $\IZ X$ is used.}
 Indeed, the hypothesis on 
$\cA_X=\pi^{-1}(S)$ of the reference is satisfied since $\IZ X$ is
Zariski dense in $\cA$. 
We get a quotient abelian
scheme $\varphi \colon \cA \rightarrow \cB'$ of relative dimension
$g'$, \textit{i.e.}, there exists an abelian subscheme $\cB$ of $\cA
\rightarrow S$ with $\varphi$ being the quotient $\cA \rightarrow
\cA/\cB$, such that for the diagram
% induced by the quotient $\varphi \colon \cA \rightarrow \cA/\cB$
%=======
% Let $X \subseteq \cA$ be an irreducible subvariety satisfying the assumptions of Theorem~\ref{MainThm}. If $\dim X = 0$ then we are done. Assume from now on $\dim X \ge 1$. 
% By Theorem~\ref{TheoremFirstDegLocusIsX}, $X^{\mathrm{deg}}(1)$ is
% Zariski dense in $X$. So $X^{\mathrm{deg}}(1) = X$ by the Zariski
% closedness of $X^{\mathrm{deg}}(1)$,
% Theorem~\ref{ThmDegLocusZarClosed}. Now we apply
% \cite[Thm.8.1]{GaoBettiRank} with $t = 1$ and get a quotient abelian
% scheme $\varphi \colon \cA \rightarrow \cB$ of relative dimension
% $g'$, \textit{i.e.} there exists an abelian subscheme $\cB'$ of $\cA
% \rightarrow S$ such that $\varphi$ is the quotient $\cA \rightarrow
% \cA/\cB'$, and a polarization type $D'$, such that for the diagram
%>>>>>>> parent of fddc4d3 (Added a section for the criterion in terms of Betti rank. Added comments in Introduction)
%To apply APPENDIX with $t = 1$ we furthermore need the Zariski closedness of $X^{\mathrm{deg}}(1)$, Theorem~\ref{ThmDegLocusZarClosed} with $t=1$, to conclude that $X^{\mathrm{deg}}(1) = X$. Hence by APPENDIX with $t = 1$, we obtain 
    \begin{equation*}
      \xymatrix{
        \cA \ar[d]\ar[r]^{\varphi}&      \cB' \ar[d] \ar[r]^-{\iota} \pullbackcorner & \mathfrak{A}_{g'}\ar[d] \\
        S \ar@{=}[r]&      S \ar[r]^-{\iota_S} & \mathbb{A}_{g'},
      }
    \end{equation*}
we have $\dim (\iota\circ\varphi)(X) < \dim X - (g-g') + 1$ and that $\iota\circ \varphi$ is not generically finite.

We have $0 \le \dim (\iota\circ\varphi)(X) \le \dim X - (g-g')$. So if
$g'=0$ then $\dim X \ge g$ and we are done. Let us now assume $g'\ge
1$. 

We claim that $g' < g$. Indeed, otherwise $g' = g$ and $\varphi$ is the identity. But then $\iota$ is the inclusion by {\tt (Hyp)}. This contradicts the fact that $\iota\circ \varphi$ is not generically finite.% Here each vertical morphism is an abelian scheme, $\varphi$ is surjective, and the right box is Cartesian.

Now $\iota(\cB') \rightarrow \iota_S(S)$ is an abelian scheme of
relative dimension $g' \le g-1$, and $(\iota\circ\varphi)(X)$ is an
irreducible subvariety of $\iota(\cB')$.
Recall that $\IZ X$ is Zariski dense in $\cA$. It is not hard to check that $(\iota\circ\varphi)(X)$ dominates $\iota_S(S)$, $\mathbb{Z}(\iota\circ\varphi)(X)$ is Zariski dense in $\iota(\cB')$, and $(\iota\circ\varphi)(X) \cap \iota(\cB')_{\mathrm{tor}}$ is Zariski dense in $(\iota\circ\varphi)(X)$. In other words,  $(\iota\circ\varphi)(X) \subseteq \iota(\cB') \rightarrow \iota_S(S)$ still satisfy the assumptions of Theorem~\ref{MainThm}.

We are almost ready to apply the induction hypothesis, except that $\iota_S(S)$ may not be regular. Set $S_0 := \iota_S(S)^{\mathrm{reg}}$, $\cB'_0 := \iota(\cB') \times_{\iota_S(S)} S_0$ and $X_0 := (\iota\circ\varphi)(X) \cap \cB'_0$. Then $X_0$ is Zariski open dense in $(\iota\circ\varphi)(X)$, and $X_0 \subseteq \cB'_0 \rightarrow S_0$ still satisfy the assumptions of Theorem~\ref{MainThm}.

The relative dimension of $\cB'_0 \rightarrow S_0$ is $g' \le g-1$. So we can apply the induction hypothesis and get $\dim X_0 \ge g'$. So $\dim X > \dim (\iota\circ\varphi)(X) + g-g'-1 = \dim X_0 + g-g'-1 \ge g-1$. Therefore $\dim X \ge g$ and we are done.
\end{proof}

%%% Local Variables:
%%% TeX-master: "main"
%%% End:


%% Section 8
\section{Uniformity of the Number of Algebraic Torsion Points in
  Curves}
\label{SectionUnifMM}
In this section, we explain how Theorem~\ref{MainThm} for $L=\IQbar$
implies Corollary~\ref{CorUniformMM}.


Let $\mathbb{M}_g$ be the moduli space of smooth projective
geometrically irreducible curves of genus $g\ge 2$ with symplectic
level-$3$-structure. Let $\mathfrak{C}_g \rightarrow \mathbb{M}_g$
be the universal curve. Let $\mathfrak{J}_g =
\mathrm{Jac}(\mathfrak{C}_g/\mathbb{M}_g) = \mathrm{Pic}^0(\mathfrak{C}_g/\mathbb{M}_g)$ be the relative Jacobian.
As usual, everything is in characteristic $0$ and we take
$\mathfrak{C}_g$ to be defined over a number field and geometrically
irreducible.

For the proof of Corollary~\ref{CorUniformMM} we treat familes of
curves over $\mathbb{M}_g$. The proof of this
lemma is inspired by \cite[Theorem~2.4]{Stoll:Uniform} and \cite[Proposition~7.1]{DGHUnifML}. For each subvariety $S$ of $\mathbb{M}_g$, write
$\mathfrak{C}_S$ and $\mathfrak{J}_S$ for the base changes
$\mathfrak{C}_g\times_{\mathbb{M}_g}S$ and
$\mathfrak{J}_g\times_{\mathbb{M}_g}S$.

\begin{lemma}\label{LemUML}
  Let $S$ be a regular, irreducible, locally closed subvariety of
  $\mathbb{M}_g$. Then there exists a constant $c = c(S) > 0$ such that
  for all $s\in S(\IQbar)$ there is $\Xi_s\subset
  \mathfrak{C}_s(\IQbar)$ with $\#\Xi_s < c$ and with the following
  property.
  For all $x\in \mathfrak{C}_s(\IQbar)\ssm \Xi_s$ we have
  % \begin{enumerate}
  % \item[(i)] Either $x$ belongs to a subset of $\mathfrak{C}_s(\IQbar)$ of cardinality at most $c$,
  % \item[(ii)] or
    $\#\{y \in \mathfrak{C}_s(\IQbar) : y-x \in (\mathfrak{J}_s)_{\mathrm{tor}} \} < c$.
%  \end{enumerate}
\end{lemma}
%Here $Q-P$ in (ii) is defined as follows: let $\mathfrak{C}_s - P$ be the image of the Abel--Jacobi map of $\mathfrak{C}_s$ to $\mathfrak{J}_s$ based at $P$.

\begin{proof}
We prove the lemma by induction on $\dim S$. The proof for the base step $\dim S =0$ is in fact contained in the induction.

Alternatively, the base case $\dim S =0$ follows from work of Baker
and Poonen as follows. If $S$ is a point $s$, we are dealing with a single
curve. We may take $\Xi_s=\emptyset$. The desired bound follows from
\cite[Corollary~3]{BakerPoonen:01} on torsion packets.

Set $ \mathfrak{C}_g^{[6]} $ to be the $6$-th fibered power of
$\mathfrak{C}_g$ over $\mathbb{M}_g$, and $\mathfrak{J}_g^{[5]}$ to be
the $5$-th fibered power of $\mathfrak{J}_g$ over $\mathbb{M}_g$. Let
$\mathscr{D}_5 \colon \mathfrak{C}_g^{[6]} \rightarrow
\mathfrak{J}_g^{[5]}$ be the Faltings--Zhang map, fiberwise defined by
sending $(x_0,x_1,\ldots,x_5) \mapsto (x_1-x_0,\ldots,x_5-x_0)$. Write
$\mathfrak{C}_S^{[6]}$ and $\mathfrak{J}_S^{[5]}$ for the base changes
$\mathfrak{C}_g^{[6]}\times_{\mathbb{M}_g}S$ and
$\mathfrak{J}_g^{[5]}\times_{\mathbb{M}_g}S$.


Consider $X= \mathscr{D}_5(\mathfrak{C}_S^{[6]})$, which is an
irreducible subvariety of $\mathfrak{J}_S^{[5]}$.

Observe that $\dim X\le 6+\dim S$. Since $g \ge 2$, we have $$ \dim
{\mathfrak{J}_S^{[5]}} - \dim X \ge 5g-6 > 3g-3 = \dim
\mathbb{M}_g\ge\dim S.$$
We conclude $\dim X < \dim \mathfrak{J}_S^{[5]}-\dim S$.


We claim that $\ZZ X = \bigcup_{N \in \ZZ} [N]X$ is Zariski dense in
$\mathfrak{J}_S^{[5]}$. Indeed, it suffices to prove that for each $s
\in S(\IQbar)$, the fiber $X_s$ is not contained in any proper
algebraic subgroup of $\mathfrak{J}_s^5$. Now take  $x_0 \in \mathfrak{C}_s(\IQbar)$, then $X_s = \mathscr{D}_5(\mathfrak{C}_s^6) \supseteq \mathscr{D}_5(\{x_0\} \times \mathfrak{C}_s^5) = (\mathfrak{C}_s-x_0)^5$, with $\mathfrak{C}_s - x_0$ viewed as the image of the Abel--Jacobi map from $\mathfrak{C}_s$ to $\mathfrak{J}_s$ based at $x_0$. Since $\mathfrak{C}_s - x_0$ generates $\mathfrak{J}_s$, we conclude that $X_s$ is not contained in any proper subgroup of $\mathfrak{J}_s^5$.


%The geometric generic fiber of $X$ over $S$ is irreducible. %It is not
%a proper algebraic subgroup of the geometric generic fiber of
%$\mathfrak{J}_S^{[5]}\rightarrow S$. !!!TODO: Explain


Thus we can apply Theorem~\ref{MainThm} to $X$  in
$\mathfrak{J}_S^{[5]} \rightarrow S$ and $L=\IQbar$. We conclude that 
\[
Y = \overline{X \cap (\mathfrak{J}_S^{[5]})_{\mathrm{tor}}}^{\mathrm{Zar}}
\]
is a proper, Zariski closed subset of $X$. 



Let $\pi \colon \mathfrak{J}_S^{[5]} \rightarrow S$ denote the
structure morphism.
Now $\dim \overline{S \setminus
  \pi(X\setminus Y)}^{\mathrm{Zar}}< \dim S$.
Thus $\overline{S \setminus
  \pi(X\setminus Y)}^{\mathrm{Zar}}$ is a finite union of regular,
irreducible, locally closed subvarieties $S'$ of $S$ with $\dim
S'<\dim S$. 
So we can apply induction hypothesis to each of these $S'$ and
conclude the result. Thus, in order to prove the
lemma, it suffices to consider $s\in S(\IQbar)$  in
$\pi(X\setminus Y)$. Thus we have % we may replace $S$ by $S \setminus \bigcup S_i$, and therefore assume $\pi(X \setminus X_{\mathrm{at}}) = S$. So
\begin{equation*}
Y_s \subsetneq X_s = \mathscr{D}_5(\mathfrak{C}_s^{[6]}). % \quad \text{ for all } s \in S(\IQbar).
\end{equation*}
Let $x \in \mathfrak{C}_s(\IQbar)$.
% Taking into consideration this
% inclusion, 
We are in one of the two cases:
\begin{enumerate}
\item[(i)] $(\mathfrak{C}_s-x)^5 \subseteq Y_s$ or % \subsetneq \mathscr{D}_5(\mathfrak{C}_s^6)$;
\item[(ii)] $(\mathfrak{C}_s-x)^5 \not\subseteq Y_s$.
\end{enumerate}

In case (i), we apply \cite[Lemma 6.4]{DGHUnifML} with
$Z$ any irreducible component of $Y_s\subsetneq
\mathscr{D}_5(\mathfrak{C}_s^6)$. We conclude that
there are at most $84(g-1)\cdot n(Y_s)$ possibilities
for $x$ where $n(Y_s)$ is the number of irreducible
components of $Y_s$. Notice that
$n(Y_s)$ is strictly less than a number $c_1'$ independent of $s$. Set $c_1 = 84(g-1)c_1'$.
The $x$ leading to case (i) will amount to the points in $\Xi_s$. %from case (i) will be in alternative (i) from this proposition.

In case (ii), set $\Sigma = \{y \in \mathfrak{C}_s(\IQbar) : y-x \in
(\mathfrak{J}_s)_{\mathrm{tor}} \}$. By \cite[Lemma 6.3]{DGHUnifML},
there exists a number $c_2$, independent of $s$, such that the
following holds. If $\# \Sigma  \ge c_2$, then $(\Sigma-x)^5 \not\subseteq Y_s(\IQbar)$. But $(\Sigma-x)^5 = \mathscr{D}_5(\{x\} \times \Sigma^5) \subseteq X(\IQbar)$ and $(\Sigma-x)^5 \subseteq (\mathfrak{J}_s^{[5]})_{\mathrm{tor}}$, so $(\Sigma-x)^5 \subseteq Y_s(\IQbar)$ by definition of $Y$. Therefore $\#\Sigma < c_2$.

Now the lemma follows by taking $c = \max\{c_1,c_2\}$.
\end{proof}

Lemma~\ref{LemUML} improves itself as follows.

\begin{proposition}\label{PropUML}
  Let $S$ and $c=c(S)$ be as in Lemma~\ref{LemUML}. For
  all $s\in S(\IQbar)$ and all $x \in \mathfrak{C}_s(\IQbar)$ we
  have  $\#\{y \in \mathfrak{C}_s(\IQbar) : y-x \in (\mathfrak{J}_s)_{\mathrm{tor}} \} < c$.
\end{proposition}
\begin{proof}
  Let $s\in S(\IQbar)$ and $x\in \mathfrak{C}_s(\IQbar)$.
  Suppose $y_1,\ldots,y_n\in\mathfrak{C}_s(\IQbar)$ are pairwise
  distinct with $y_i-x$ torsion for all $i$.
  If $n\ge c$, then some $y_i$, say $y_1$, lies in
  $\mathfrak{C}_s(\IQbar)\ssm \Xi_s$ with $\Xi_s$ the set from
  Lemma~\ref{LemUML}.
  Now $y_i-y_1 = (y_i-x)-(y_1-x)$ are $n$  torsion points
  of $\mathfrak{J}_s(\IQbar)$ for $i\in\{1,\ldots,n\}$. So
  Lemma~\ref{LemUML} implies $n<c$, a contradiction. We conclude $n<c$. 
\end{proof}


\begin{proof}[Proof of Corollary~\ref{CorUniformMM}]
  By a specialization argument \cite[Lemma~3.1]{DGHBog} based on
  Masser's \cite{masser1989specializations}, it suffices to prove the
  result with $\IC$ replaced by $\IQbar$. Then the result follows
  immediately from Propsition~\ref{PropUML} applied to $S =
  \mathbb{M}_g$.
\end{proof}


%%% Local Variables:
%%% TeX-master: "main"
%%% End:



%% Section 9
\section{A Criterion for Zariski density over $\IQbar$}
\label{SectionCriterionTorsionDenseOverIQbar}
We prove Theorem~\ref{ThmCritTorsionDense} for $L=\IQbar$ in this
section. This proof uses Theorem~\ref{MainThm} for $L=\IQbar$, which
has already been establised in $\mathsection$\ref{SectionConclusion},
and the criterion of $\mathrm{rank}_{\mathrm{Betti}}(X) < 2g$ by
\cite[Theorem~1.1]{GaoBettiRank} (with $l = g$).
Recall that $\cA$ is an abelian scheme of relative dimension $g\ge 1$
over a regular, irreducible, quasi-projective base $S$ defined over $\IQbar$. 
% The proof goes as follows. We start by proving Theorem~\ref{ThmCritTorsionDense} with $\IC$ replaced by $\IQbar$. This step uses Theorem~\ref{MainThm} over $\IQbar$ and the criterion of  by \cite[Thm.1.1]{GaoBettiRank} (with $l = g$). Then we use a simple specialization argument to extend the result over $\IC$. 

%And then we can conclude the full Theorem~\ref{MainThm} over $\IC$ because it follows immediately from Theorem~\ref{ThmCritTorsionDense} and the naive bound on the generic Betti rank \eqref{EqBettiRank}.

\begin{proof}[Proof of Theorem~\ref{ThmCritTorsionDense} for $L=\IQbar$]
The implication ``$\Leftarrow$'' follows from
\cite[Proposition~2.1.1]{ACZBetti} for subfields of $\IC$.
%% Justified why the points are Qbar points and not just CC points.
Indeed,
$X(\IC)\cap \cA_{\mathrm{tors}}$ is dense in $X^{\mathrm{an}}$.
Among these point we may find a Zariski dense subset of
$\IQbar$-points as follows. 
The
Betti map is locally injective. So we may find a Zariski dense set of
isolated intersection points in $X(\IC)\cap \ker[n]$ as $n\in\IN$
varies.
These points are defined over $\IQbar$ as $X$ is. 

From now on, we focus on ``$\Rightarrow$''. Let $\cA \rightarrow S$ be an abelian scheme defined over $L=\IQbar$, and let $X$ be an irreducible subvariety of $\cA$ defined over $L$ such that $\ZZ X$ is Zariski dense in $\cA$ and that $X(L) \cap \cA_{\mathrm{tor}}$ is Zariski dense in $X$. 
(Observe that Theorem~\ref{MainThm} in the already proved case
$L=\IQbar$ now implies $\dim X\ge g$.)

Assume $\mathrm{rank}_{\mathrm{Betti}}(X) < 2g$. We wish to get a contradiction.

By \cite[Theorem~1.1]{GaoBettiRank} applied to $l = g$, there exists a quotient abelian scheme $\varphi \colon \cA \rightarrow \cB'$ of relative dimension $g'$, \textit{i.e.}, there exists an abelian subscheme $\cB$ of $\cA \rightarrow S$ with $\varphi$ being the quotient $\cA \rightarrow \cA/\cB$, such that for the diagram
    \begin{equation*}
      \xymatrix{
        \cA \ar[d]\ar[r]^{\varphi}&      \cB' \ar[d] \ar[r]^-{\iota} \pullbackcorner & \mathfrak{A}_{g'}\ar[d] \\
        S \ar@{=}[r]&      S \ar[r]^-{\iota_S} & \mathbb{A}_{g'},
      }
    \end{equation*}
we have $\dim (\iota\circ\varphi)(X) < g'$. 

Set $S_0 := \iota_S(S)^{\mathrm{reg}}$, $\cA_0' := \mathfrak{A}_{g'} \times_{\mathbb{A}_{g'}} S_0$ and $X_0 := (\iota\circ \varphi)(X) \cap \cA_0'$. Then $\cA_0' \rightarrow S_0$ is an abelian scheme of relative dimension $g'$ defined over $L$, and $X_0$ is Zariski open dense in $(\iota\circ\varphi)(X)$. So $\dim X_0 < g'$.

It is not hard to check that $\ZZ X_0$ is Zariski dense in $\cA_0'$
since $\ZZ X$ is Zariski dense in $\cA$, and that $X_0(L) \cap
\cA'_{0,\mathrm{tor}}$ is Zariski dense in $X_0$ since $X(L) \cap
\cA_{\mathrm{tor}}$ is Zariski dense in $X$. Thus we can apply
Theorem~\ref{MainThm}, over the base field $L$,
to $X_0 \subseteq \cA'_0 \rightarrow S_0$ and conclude that $\dim X_0 \ge g'$.

The conclusions of the last two paragraphs are contradictory. Thus we get the desired contradiction. 
So $\mathrm{rank}_{\mathrm{Betti}}(X) = 2g$. We are done.% with $\IC$ replaced by $\IQbar$.
\end{proof}

Let $L$ be an algebraically closed subfield of $\IC$.
The argument above shows that   Theorem \ref{MainThm} for $L$ implies Theorem
\ref{ThmCritTorsionDense} for $L$.


%%% Local Variables:
%%% TeX-master: "main"
%%% End:


%% Section 10
\section{Specialization: From $\IQbar$ to $\IC$}
\label{SectionSpecialization}
The goal of this section is to prove Theorem~\ref{MainThm}, proved in
$\mathsection$\ref{SectionConclusion} for $L=\IQbar$, for an arbitrary
algebraically closed field $L$ of characteristic $0$. Note first, that
it suffices to consider only subfields $L$ of $\IC$ by a suitable
Lefschetz principle. Moreover, as all varieties are defined using
finitely many polynomials and coefficients, we may assume that
$\mathrm{trdeg}_\IQbar L <\infty$. 

In this section we will also complete the proof of Theorem~\ref{ThmCritTorsionDense}.

Let $L$ be an algebraically closed subfield of $\IC$ that has finite
transcendence degree over $\IQbar$. Let $\cA
\rightarrow S$ be an abelian scheme defined over $L$, and let $X
\subseteq \cA$ be an irreducible subvariety. We suppose that $\ZZ X$
is Zariski dense in $\cA$ and $X(\IC) \cap \cA_{\mathrm{tor}}$ is
Zariski dense in $X$.

We shall prove both Theorem~\ref{MainThm} and
Theorem~\ref{ThmCritTorsionDense} by induction on $
\mathrm{trdeg}_{\IQbar} L$. We shall proceed as follows. For each
integer $d \ge 0$, we set
\begin{center}
%  \label{def:RMMd}
  {\tt RMM(d):} Theorem~\ref{MainThm} holds true if
  $\mathrm{trdeg}_\IQbar L\le d$ % when we replaced $\IC$ by any algebraically closed field of transcendence degree $\le d$ over $\IQbar$
\end{center}
and
\begin{center}
%  \label{def:TorDensed}
  {\tt TorDense(d):} Theorem~\ref{ThmCritTorsionDense} holds true if
  $\mathrm{trdeg}_\IQbar L\le d$. %holds true  when we replaced $\IC$ by any algebraically closed field of transcendence degree $\le d$ over $\IQbar$.
\end{center}
We proceed by proving the following statements
for each integer $d \ge 0$:
\begin{enumerate}
\item[(i)] {\tt RMM(d)} implies {\tt TorDense(d)};
\item[(ii)] {\tt TorDense(d)} implies {\tt RMM(d+1)}.
\end{enumerate}
Note that the
statement ${\tt RMM(0)}$ was proved in
$\mathsection$\ref{SectionConclusion}, \textit{i.e.},
Theorem~\ref{MainThm} in the case $L=\IQbar$.
% Then we can conclude that both Theorem~\ref{MainThm} and
% Theorem~\ref{ThmCritTorsionDense} because we have proved {\tt RMM(0)}
% in 


\subsection{Proof of {\tt RMM(d)}$\Rightarrow${\tt TorDense(d)}}
This follows from a verbalized copy of the proof executed in $\mathsection$\ref{SectionCriterionTorsionDenseOverIQbar}.%and notice that we have proved the full version of Theorem~\ref{MainThm} (over $\IC$, not only over $\IQbar$), we can conclude for Theorem~\ref{ThmCritTorsionDense}.





\subsection{Proof of {\tt TorDense(d)}$\Rightarrow${\tt RMM(d+1)}}
%The direction $\Leftarrow$ is proved by \cite[Prop.2.1.1]{ACZBetti}. From now on, we focus on $\Rightarrow$. 
%The base step $d = 0$ is the case where all varieties in question are defined over $\IQbar$, and is proved in $\mathsection$\ref{SectionConclusion}.
%Now let $d \ge 1$. 
%Assume that Theorem~\ref{MainThm} is proved for $0,\ldots,d-1$. We wish to prove it for $d$. 
Suppose we are in the case $\mathrm{trdeg}_{\IQbar} L = d+1$. So
$\cA,S,$ and $X$ are all defined over $L$. We assume that $\IZ X$ is
Zariski dense in $\cA$ and
$X(L)\cap\cA_{\mathrm{tors}}$ is Zariski dense in $X$. In particular,
$S=\pi(X)$.
Our goal is to
show $\dim X \ge g$ with $g$ the relative dimension of $\cA/S$.

To start we follow the argument from the beginning of the proof of
Theorem~\ref{MainThm} to reduce to the universal family.
We obtain from $X$ a regular, irreducible, locally closed $S''\subset\IA_{g,L}$
as well as an irreducible subvariety  $X''\subset
\cA''=\mathfrak{A}_{g,L}\times_{\IA_{g,L}} S''$
such that $\IZ X''$ is Zariski dense in $\cA''$
and such that $X''(L)\cap\cA''_{\mathrm{tors}}$ is Zariski dense in
$X''$. Moreover, $\dim X'' \le \dim X$. If we can establish $\dim
X''\ge g$ then we are done.

From now on we assume that $X$ is an irreducible  subvariety of
$\mathfrak{A}_{g,L}$ and $\pi(X) = S\subset\IA_{g,L}$. Recall that we
already reduced to the case $L\subset\IC$.

There exists an algebraically closed subfield $K \subseteq L$ with 
 $K\supset \IQbar$ such that $\mathrm{trdeg}_{\IQbar}K = d$. Then
$\mathrm{trdeg}_K L = 1$.

By \cite[Lemma~2.2]{BarDill} there is an irreducible  subvariety
$\mathfrak{X}\subset \mathfrak{A}_{g,K}$ such that $X\subset
\mathfrak{X}_L$ and $\dim \mathfrak X \le \dim X + 1$. Note that $\dim
X\le \dim\mathfrak X$. We may assume that $\mathfrak X$ is the
minimal subvariety of $\mathfrak{A}_{g,K}$ whose base change to $L$
contains $X$.

If $\dim \mathfrak X =\dim X$, then $X$ was originally already
defined over $K$. In this case $\dim X\ge g$ follows from \texttt{RMM(d)}.
So we may assume $\dim \mathfrak X = \dim X + 1$. 

Let $\mathfrak S = \pi(\mathfrak X)\subset\IA_{g,K}$.
Then $\mathfrak{S}_L\supset S$.
We claim that $\IZ \mathfrak X$ is Zariski dense in
$\pi^{-1}(\mathfrak S)$.
The Zariski closure $Z$ of $\IZ \mathfrak X$ in
$\pi^{-1}(\mathfrak S)$ %$\cA$ %$\mathfrak{A}_{g,K}$
has dimension at least
$g+\dim S = \dim \pi^{-1}(S)$ since $Z_L\supset \overline{\IZ
  X}^{\mathrm{Zar}}$.
Our claim follows if $\dim S =
\dim\mathfrak{S}$.  So we may assume $1+\dim S \le  \dim \mathfrak S$.
Suppose $g+\dim S = \dim Z$, then  $Z_L= \pi^{-1}(S)$
which contradicts $ \pi(Z)=\mathfrak{S}$.
 Thus $\dim Z\ge g + \dim S + 1$.
By \cite[Lemma~2.2]{BarDill} $S$ is contained in $\mathfrak{S}'_L$
with $\mathfrak{S}'\subset\IA_{g,K}$ and $\dim \mathfrak{S}' \le \dim
S + 1$. So $\mathfrak{X} \subset\pi^{-1}(\mathfrak S')$ by minimality
of $\mathfrak{X}$. Thus $\mathfrak S \subset \mathfrak S'$ and in
particular $\dim \mathfrak{S} = \dim S + 1$.
So we must have $\dim Z \ge
g+\dim S + 1 = g+\dim\mathfrak S = \dim \pi^{-1}(\mathfrak S)$. This
implies $Z=\pi^{-1}(\mathfrak S)$.
%% Not necessary:
% But 
% $\dim Z\le \dim \pi^{-1}(\mathfrak S) = g+\dim\mathfrak S$.
%%% Old version
% Our claim follows if $\dim S =
% \dim\mathfrak{S}$.  Otherwise we have $\dim S < \dim \mathfrak S$.
% By \cite[Lemma~2.2]{BarDill} $S$ is contained in $\mathfrak{S}'_L$
% with $\mathfrak{S}'\subset\IA_{g,K}$ and $\dim \mathfrak{S}' \le \dim
% S + 1$. So $\mathfrak{X} \subset\pi^{-1}(\mathfrak S')$ by minimality
% of $\mathfrak{X}$. Thus $\mathfrak S \subset \mathfrak S'$ and in
% particular $\dim \mathfrak{S} = \dim S + 1$. Thus $g+\dim S \le \dim
% Z\le g + \dim S + 1$. If $g+\dim S = \dim Z$, then $Z_L= \pi^{-1}(S)$
% which contradicts $ \pi(Z)=\mathfrak{S}$. So we must have $\dim Z =
% g+\dim S + 1 = g+\dim\mathfrak S = \dim \pi^{-1}(\mathfrak S)$. This
% implies $Z=\pi^{-1}(\mathfrak S)$.
%%%

We claim that $\mathfrak{X}_\IC^{\mathrm{deg}}(0)$ lies Zariski dense
in $\mathfrak{X}$.
Let $x \in X(L) \cap (\mathfrak{A}_{g,L})_{\mathrm{tors}}$.
The $K$-Zariski closure $C$ of $\{x\}$ in $\mathfrak{X}$ is contained in the kernel of
$[\mathrm{ord}(x)]$.
We have $\dim C\le 1$ since $\mathrm{trdeg}_K L = 1$, indeed, use
again \cite[Lemma~2.2]{BarDill}. As 
$X(L) \cap (\mathfrak{A}_{g,L})_{\mathrm{tors}}$ is Zariski dense in
$X$ and as $X$ is not defined over $K$, we may assume $\dim C = 1$. 
Such a curve lies in $\mathfrak{X}_\IC^{\mathrm{deg}}(0)$ as it is torsion. Ranging
over the possible $x$ yields our claim.

By \cite[Theorem~1.7]{GaoBettiRank} we conclude
$\mathrm{rank}_{\mathrm{Betti}}(\mathfrak{X}) < 2\dim \mathfrak{X}$.

The above argument also implies that
$\mathfrak{X}(K)\cap(\mathfrak{A}_{g,K})_{\mathrm{tors}}$ is Zariski
dense in $\mathfrak X$. We apply \texttt{TorDense(d)} to
$\mathfrak{X}\cap\pi^{-1}(\mathfrak{S}^{\mathrm{reg}})$, which is
defined over $K$, and conclude
$\mathrm{rank}_{\mathrm{Betti}}(\mathfrak{X}) =2g $.

Combining both bounds for the Betti rank yields $\dim\mathfrak X > g$.
As $\dim X + 1 = \dim \mathfrak{X}$ we conclude $\dim X\ge g$, as desired. \qed
\iffalse

There exists a regular, irreducible curve $C$ defined over $K$ with the following properties: all varieties in question $X \subseteq \cA \rightarrow S$ are defined over the function $K(C)$, and $L$ is an algebraic closure of $K(C)$. Up to replacing $C$ by a Zariski open dense subset, we may assume that there exists a $C$-scheme $\mathfrak{S} \rightarrow C$ such that the geometric generic fiber is $S$. Notice that $\mathfrak{S}$ is a variety defined over $K$.

Let us spread $\cA \rightarrow S$ to an abelian scheme over $\mathfrak{S}$ as follows. There exists an $S$-closed immersion $\cA \rightarrow \IP_S^N$ for some $N > 0$. Thus we have a morphism $\cA \rightarrow \IP_S^N \rightarrow \IP_{\mathfrak{S}}^N$. Let $\mathfrak{A}$ be the closure of the image of this map. Up to replacing $C$ by a Zariski open dense subset (and hence $\mathfrak{S}$ by a Zariski open dense subset), the $S$-group structure on $\cA$ extends to an $\mathfrak{S}$-group structure on $\mathfrak{A}$. Thus $\mathfrak{A} \rightarrow \mathfrak{S}$ is an abelian scheme, because it is a projective group scheme. Notice that $\cA = \mathfrak{A} \times_{\mathfrak{S}} S$ by construction, and $\mathfrak{A} \rightarrow \mathfrak{S}$ is defined over $K$.

To sum it up, we have the following Cartesian diagram% with $\iota'$ the modular map,
\begin{equation*}
  \xymatrix{
    \cA \ar[r]^-{i} \ar[d]_-{\pi} \pullbackcorner & \mathfrak{A} \ar[d]^-{\pi_{\mathfrak{S}}} \\
    S \ar[r]^-{i_S}   & \mathfrak{S}.
  }
\end{equation*}
Let $\mathfrak{X}$ be the Zariski closure of $i(X)$. 
Then $\mathfrak{X}$ is an irreducible subvariety of $\mathfrak{A}$ defined over $K$, and $\ZZ \mathfrak{X}$ is Zariski dense in $\mathfrak{A}$ because $\ZZ X$ is Zariski dense in $\cA$. Moreover, we have $\dim \mathfrak{X} \le \dim X + 1$ by \cite[Lemma~2.2]{BarDill}: $K \subseteq L$ is an extension of algebraically closed fields of transcendence degree $1$, $\mathfrak{A}$ is a variety defined over $K$ and $X$ is a subvariety of $\mathfrak{A}\otimes_K L$. 






We claim that: (i) $\mathfrak{X}(K) \cap \mathfrak{A}_{\mathrm{tor}}$ is Zariski dense in $\mathfrak{X}$; (ii) $\mathfrak{X}^{\mathrm{deg}}(0)$ is Zariski dense in $\mathfrak{X}$. 
 Indeed, since $X\subseteq \cA \rightarrow S$ are defined over $L$ and $X(\IC) \cap \cA_{\mathrm{tor}}$ is Zariski dense in $X$, we have that $X(L) \cap \cA_{\mathrm{tor}}$ is Zariski dense in $X$. For each point $x \in \tor{\cA(L)}$, let $\mathfrak{S}_x$ be the Zariski closure of $i_S(\pi(x))$. Then $\mathfrak{S}_x$ is an irreducible subvariety of $\mathfrak{S}$ defined over $K$ of dimension $1$, and the Zariski closure of $i(x)$ is a torsion multisection of $\mathfrak{A}\times_{\mathfrak{S}} \mathfrak{S}_x \rightarrow \mathfrak{S}_x$. 
Therefore, $\mathfrak{X}(K) \cap \mathfrak{A}_{\mathrm{tor}}$ is Zariski dense in $\mathfrak{X}$ because $X(L) \cap \cA_{\mathrm{tor}}$ is Zariski dense in $X$. On the other hand, $i(x) \subseteq \mathfrak{X}^{\mathrm{deg}}(0)$ by definition of the degeneracy locus. So $\mathfrak{X}^{\mathrm{deg}}(0)$ is Zariski dense in $\mathfrak{X}$ because $X(L)\cap \cA_{\mathrm{tor}}$ is Zariski dense in $X$.


By (i), we can apply {\tt TorDense(d)} to $\mathfrak{X} \subseteq \mathfrak{A} \rightarrow \mathfrak{S}$ (which are defined over $K$ of transcendence degree $d$ over $\IQbar$). We thus obtain
\[
\mathrm{rank}_{\mathrm{Betti}}(\mathfrak{X}) = 2g.
\]
By (ii) we can apply 
\cite[Theorem~4.3.1]{GaoHDR} to conclude that
$\mathrm{rank}_{\mathrm{Betti}}(\mathfrak{X}) < 2\dim \mathfrak{X}$.
Thus $g<\dim\mathfrak X$. As $\dim\mathfrak X \le 1+\dim X$ we find
$\dim X\le g$, which is what we desire.
\fi

%%% Local Variables:
%%% TeX-master: "main"
%%% End:




\appendix
\renewcommand{\thesection}{\Alph{section}}
\setcounter{section}{0}

\section{Large Galois Orbits revisited}
\label{Appendix}
In $\mathsection$\ref{SectionLGO} we used a quantitative version,
obtained by R\'{e}mond \cite[Proposition~2.9]{RemondGaloisBound}, of
Masser's earlier result on the Galois orbit of torsion points on
abelian varieties. In this appendix, we prove in
Proposition~\ref{PropDavidBound} an estimate that is sufficient for
our purposes. We will rely on a result of David
\cite[Th\'eor\`eme~1.4]{DavidMinHaut}, it is not directly related to
Masser and W\"{u}stholz's Isogeny Theorem \cite{MW:abelianisog}.


Throughout this section let $k$ be a number field contained in a fixed algebraic closure
$\overline k$. 
If not stated otherwise,  $A$ denotes
an abelian variety of dimension $g\ge 1$  defined over  $k$. We set
\[
  \rho(A,k) = 
  [k:\IQ](\max\{1,h_{\mathrm{Fal}}(A)\}+\log [k:\IQ])\ge [k:\IQ]% + D^{1/(\dim A+2)}
\]
%with $h = $ and
where
$h_{\mathrm{Fal}}(A)$ denotes the \textit{stable Faltings height} of
$A$. The additive normalization in $h_{\mathrm{Fal}}$ plays no role in
the current work.

We begin by stating a special case of \cite[Th\'eor\`eme 1.4]{DavidMinHaut}. 

\begin{theorem}[David]
  \label{thm:david}
  For each integer $g\ge 1$  there is a constant $c_1(g)>0$
  with the following property. 
  Let $(A,\cL)$ be
  a principally polarized abelian variety of dimension $g$
  defined over $k$. Suppose $x\in A(k)$
  has finite order. Then there exists an abelian subvariety
  $B\subsetneq A$ defined over $\overline k$ and  positive integer $N$ such that
  $[N](x)\in B(\overline k)$ and  $\max\{N,\deg_\cL B\} \le c_1(g) \rho(A,k)^{g+1}$.
\end{theorem}
\begin{proof}
  Since $x$ has finite order, its N\'eron--Tate height vanishes. So we
  are in the second case of \cite[Theorem~1.4]{DavidMinHaut}.
  David used a height of $A$ defined using theta functions
  (and thus is sometimes called a \textit{Theta height}). The comparison
  with  $h_{\mathrm{Fal}}(A)$ is done
  by work of Bost and David; see \cite[Corollaire~6.9]{DPvarabII} and
  \cite[Corollary~1.3]{Pazuki:12}. Thus in \cite[Th\'eor\`eme~1.4]{DavidMinHaut}, we can use the stable Faltings height after
  modifying the multiplicative constant $c_1(g)$ (which is called $c_2$
  in the reference.)
  Moreover, in the reference we  may assume $\|\mathrm{Im}\tau\|\ge
  \sqrt{3}/2$ as $\tau$ can be assumed to be Siegel reduced.
  Set $D = \max\{2,[k:\IQ]\}$ and $h = \max\{1,h_F(A)\}$. 
  Observe that
  $$
  \frac{2}{\sqrt 3} D(h+\log D) + D^{1/(g+2)} \le 2 D(h+\log D) + D\le
  3 D(h+\log D)
  $$
  as $h+\log D\ge 1$. If $k\not=\IQ$, then $D(h+\log D) = \rho(A,k)$
  and otherwise $D(h+\log D) = 2(h+\log 2) \le 2(1+\log 2) h \le 4 \rho(A,\IQ)$. 
  So the desired bound follows after
  modifying $c_1(g)$ again.
\end{proof}

We will reduce to the principally polarized case using the next lemma.

%In general we can reduce to the principally polarized case. This can be done by Masser--W\"{u}stholz's Isogeny Theorem  
% \cite{MW:abelianisog}. For our purpose we only need a weaker bound, for which we hereby give a more elementary proof. We start with the following lemma.
%using the next lemma.

\begin{lemma}
  \label{lem:ppreduction}
  Let $(A,\cL)$ be a polarized abelian
  variety with $\dim A =g\ge 1$ and
  $\Delta = (\deg_\cL A)/g!$. Then $\Delta = \dim H^0(A,\cL)$ and 
  there is a number field $k'\supset k$ 
  with $[k':k]\le \Delta^{2g}$,
  a principally polarized abelian variety $(B,\cM)$ defined over
  $k'$, and an isogeny $A_{k'}\rightarrow B$, also defined over
  $k'$, of degree at most $\Delta$. Moreover,
  $\rho(B,k')\le \Delta^{2g}(\rho(A,k) + 3g[k:\IQ]\log\Delta)$. 
\end{lemma}
\begin{proof}
  The Riemann--Roch Theorem and
  Vanishing Theorem on higher cohomology for ample line bundles,
  see \cite[Section 16]{MumfordAbVar70}, imply $\Delta=\dim H^0(A,\cL) = 
  (\deg_\cL A)/g!$. This is the first claim.
  By a classical result, see \cite[Lemme~3.5]{GR:periodsisgoenies},
  there is an isogeny $\varphi\colon A_{\overline k}\rightarrow B$ of degree
  $\dim H^0(A,\cL) =\Delta$
  to a principally polarized abelian variety $(B,\cM)$
  defined over $\overline k$; here $\varphi^* \cM = \cL$. 

  The kernel $\ker \varphi$ is isomorphic to a product of cyclic %$s$
  groups of order $t_1,\ldots,t_s$, say. Both $B$ and $\varphi$ are
  defined over the field $k'$ field generated by certain points of order
  $t_1,\ldots,t_s$, respectively. A point in $A(\overline k)$ of order
  $t_i$ generates a field extension of $k$ of degree at most $t_i^{2g}$.
  Therefore, $[k':k]\le (t_1\cdots t_s)^{2g} = \Delta^{2g}$.

  Faltings's estimate implies
  \begin{equation*}
%    \label{eq:boundhFC}
    h_F(B)\le h_F(A) + \frac 12 \log\Delta. % \le h_F(A) + \frac 32
    % \log(c_1(g)\rho(A,k)^{g+1}) +c_3(g)\le c_4(g)\frac{\rho(A,k)}{[k:\IQ]}.
  \end{equation*}
  So
  \begin{alignat*}1
    \rho(B,k') &= [k':\IQ] (\max\{1,h_F(B)\}+\log[k':\IQ]) \\
    &\le \Delta^{2g}[k:\IQ] \left(\max\{1,h_F(A)\}+\frac 12\log\Delta +
    2g\log\Delta +\log[k:\IQ]\right)\\
    &=\Delta^{2g} \rho(A,k)+\Delta^{2g}[k:\IQ]\left(\frac 12+2g\right)\log\Delta. \qquad \qquad \qedhere
  \end{alignat*}
\end{proof}


The following lemma is obtained by iterating David's result. For
an integer $g\ge 1$ we define $\lambda(g) = 3^{g+1}(g!)^2 \ge 3$. 

\begin{lemma}\label{lem:DavidBound}
  For each integer $g\ge 1$  there is a constant $c_2(g)>0$
  with the following property.   
  Let
  $(A,\cL)$ be
  a {principally polarized} abelian variety of dimension $g$
  defined over $k$. % Set $h = \max\{1,h_{\mathrm{Fal}}(A)\}$.
  % There
  % exists a constant $c_2 = c_2(g) > 0$ such that any
  If $x\in A(k)$ has finite order, then
   $\ord{x}\le c_2(g) \rho(A,k)^{\lambda(g)}$.
\end{lemma}
\begin{proof}
  The field generated by $k$ and all $3$-torsion points of $A$ has
  degree at most $3^{(2g)^2}$ over $k$. After replacing $k$  by this
  field  we may assume that all $3$-torsion points are
  $k$-rational. By Silverberg \cite[Theorem~2.4]{Silverberg:fielddef}
  and Poincar\'e's Complete Reducibility Theorem,  all abelian subvarieties of $A$ are
  defined over $k$. %%Reference: Lemma 7.2a in Masser-Wüstholz
                    %%Factorization estimates for abelian varieties (1995)

  Below $c_2(g),c_3(g),\ldots$ denote positive values that depend only
  on $g$. We will determine $c_2(g)$ by induction on $g\ge 1$.
  The case $g=1$ is treated during the induction.
  % The case $g=0$
  % is trivial, so we assume $g\ge 1$.

  Suppose $x\in A(k)$ has finite order. 
  We let $B$ and $N$ be as in  Theorem~\ref{thm:david} applied to $x$.
  If $B=0$, then $x$ has order at most $N$
  and we are done.
  In the case $g=1$ then $B=0$ since $B\subsetneq A$. So this argument
  covers the base case of the induction.
  % By the Riemann--Roch Theorem we have $\dim H^0(B,\cL|_B) = 
  % (\deg_\cL B)/(\dim B)!$.

  We assume $\dim B\ge 1$.
  Note that $B$ need not be principally polarized. Let $k'$ and $C$ be
  as in Lemma~\ref{lem:ppreduction} applied to $(B,\cL|_B)$. We define
  \begin{equation}
    \label{eq:Deltabound}
     \Delta = \dim H^0(B,\cL|_B) = \frac{\deg_\cL B}{(\dim B)!}\le c_1(g) \rho(A,k)^{g+1} 
  \end{equation}

  It is known that
  the stable Faltings height satisfies
  \begin{equation*}
    h_F(B) \le h_F(A) + \log \dim H^0(B,\cL|_B) + c_3(g) = h_F(A) +
    \log\Delta + c_3(g),
  \end{equation*}
  see \cite[$\mathsection$2.3]{GR:periodsisgoenies}. This and (\ref{eq:Deltabound}) imply
  \begin{equation*}
    \rho(B,k) \le [k:\IQ] ( \max\{1,h_F(A)\} + \log \Delta +c_3(g) +
    \log[k:\IQ])
    \le c_4(g) \rho(A,k). 
  \end{equation*}

  % By a classical result attributed to Mumford, see Lemma
  % 3.5~\cite{GR:periodsisgoenies},
  % there is an isogeny $\varphi\colon B\rightarrow C$ of degree
  % $$\Delta = \dim H^0(B,\cL|_B)\le c_1(g)\rho(A,k)^{g+1}$$
  % to a principally polarized abelian variety $C$
  % defined over $\overline k$.

  % The group $\ker \varphi$ is isomorphic to a product of cyclic $s$ groups
  % of order $t_1,\ldots,t_s$. A point in $B$ of order $t_i$ generates a
  % field extension of degree at most $t_i^{2\dim B}$. So all points in
  % $\ker \varphi$ are $k'$-rational for an extension $k'/k$ of degree
  % at most $(t_1\cdots t_s)^{2\dim B} = \Delta^{2\dim B}\le c_1(g)^{2\dim B}
  % \rho(A,k)^{2(g+1)\dim B }$. 
  %% The field
  %% $k'=k(B[\Delta])$ 
  %% generated by points in $B(\overline k)$ of order dividing
  %% $\Delta$
  %% has degree at most $\Delta^{(2\dim B)^2}\le c_1(g)^{(2\dim B)^2}
  %% \rho(A,k)^{(2\dim B)^2 (g+1)}$ over $k$.
%  Thus $C$ is defined over $k'$.

  We write $\varphi\colon B_{k'}\rightarrow C$ for the isogeny provided by
  Lemma~\ref{lem:ppreduction}. 
  As $[N](x)\in B(k')$, the image $\varphi([N](x))\in C(k')$
  is well-defined and of finite order. 
  Moreover, $\dim C=\dim B < g$, so this lemma applied by
  induction to the pair $C,k'$ and  $\varphi([N](x))$
  together with the
  estimates above
  %(\ref{eq:boundhFC})
  yields
  \begin{alignat*}1
    \ord{\varphi([N](x))} &\le  c_2(\dim B)
    \rho(C,k')^{\lambda(\dim B)} \\
    &\le c_5(g) c_2(\dim B) \left(\Delta^{2\dim B}(\rho(B,k)+[k:\IQ]
      \log\Delta)\right)^{\lambda(\dim B)}
    \\
    &\le c_6(g) c_2(\dim B)
    \left(\Delta^{2\dim B}(c_4(g)\rho(A,k)+[k:\IQ] \log\rho(A,k))\right)^{\lambda(\dim B)}
    \\
    &\le c_7(g) c_2(\dim B) \left(\Delta^{2\dim B}\rho(A,k)\right)^{\lambda(\dim B)}\\
    &\le c_8(g) c_2(\dim B) \rho(A,k)^{(1+2(g+1)\dim B)\lambda(\dim B)}. 
%     &\le c_5(g) c_2(\dim B)
%     \left([k':\IQ](\frac{\rho(A,k)}{[k:\IQ]} + \log[k':\IQ])\right)^{\lambda(\dim B)}\\
% %      \rho(A,k)^{2 (g+1)\dim B}\max\{1,h_F(C)\}\right)^{\lambda(\dim B)}\\
%     &\le c_6(g) c_2(\dim B)
%     \left(\rho(A,k)^{(1+2(g+1)\dim B)}+[k':\IQ]\log[k':\IQ]\right)^{\lambda(\dim B)}.
  \end{alignat*}
  % Now $[k':\IQ]\log [k':\IQ] \le c_7(g) [k:\IQ]\rho(A,k)^{2(g+1)\dim
  %   B} \log([k:\IQ] \rho(A,k))$. This is at most
  % $c_8(g)\rho(A,k)^{1+2(g+1)\dim B}$ by the definition of $\rho(A,k)$. We conclude
  % \begin{alignat*}1
  %   \ord{\varphi([N](x))} &\le  c_2(\dim B) c_9(g)
  %   \rho(A,k)^{(1+2(g+1)\dim B)\lambda(\dim B)}.
  % \end{alignat*}
  
  We use the bound for $N$ from Theorem~\ref{thm:david} and the bound for
  $\deg\varphi$ from Lemma~\ref{lem:ppreduction} together with
  (\ref{eq:Deltabound}) to find
  \begin{alignat*}1
    \ord{x} &\le 
    N(\deg\varphi)
    \ord{\varphi([N]x)} \\
    &\le c_1(g)^2 \rho(A,k)^{2(g+1)}
    \ord{\varphi([N]x)}
    \\
    &\le c_{9}(g) c_2(\dim B) \rho(A,k)^{2(g+1)+(1+2(g+1)\dim B)\lambda(\dim
      B)}.
%    (\log D)^{c_9(g)}.
    % D^{2(g+1)+\lambda(\dim B)+(2\dim
    %   B)^2(g+1)\lambda(\dim B)}(\log
    % D)^{c_9(g)} h^{2(g+1)+((2\dim B)^2(g+1)+1)\lambda(\dim B)}.
  \end{alignat*}

  As $\dim B\le g-1$ and $\lambda(\dim B)\le \lambda(g-1)$
the exponent of $\rho(A,k)$ is at most
  \begin{equation*}
%    2(g+1)+(1+2(g+1)\dim B)\lambda(\dim B) \le
    2(g+1)+(2g^2-1)\lambda(g-1)\le 2g +2g^2\lambda(g-1);
  \end{equation*}
  the final step used $\lambda(g-1)\ge 2$. The proposition follows as
  $2g+2g^2\lambda(g-1) = 2g+2 \cdot 3^g g^2((g-1)!)^2=2g+2\cdot 3^g
  (g!)^2 \le \lambda(g)$.
  %
% We apply \cite[Theorem 1.4]{DavidMinHaut} to the torsion point $P \in
% A(k)$ and observe that 
% Since $\hat{h}(P) = 0$, we must be in case 1 of the theorem.
% Hence there exist a constant $c_1 = c_1(g) > 0$ and an abelian
% subvariety $A_1$ of $A$ (with $\dim A_1 < g$) such that the following
% property holds true: $[M_1] P \in A_1(k)$ for some positive integer
% $M_1 \le c_1 \rho(A,k)^g (\log \rho(A,k))^g$.
%
% Apply \cite[Theorem 1.4]{DavidMinHaut} to $[M_1]P \in A_1(k)$. Then there exist a constant $c_2 = c_2(g) > 0$ and an abelian subvariety $A_2$ of $A_1$ (with $\dim A_2 < \dim A_1$)  such that $[M_2][M_1]P \in A_2(k)$ for some positive integer $M_2 \le c_2 \rho(A_1,k)^{\dim A_1} (\log \rho(A_1,k))^{\dim A_1}$. But $h_{\mathrm{Fal}}(A_1) \le h_{\mathrm{Fal}}(A)$, and hence $\rho(A_1,k) \le \rho(A,k)$. So $M_2 \le c_2 \rho(A,k)^g (\log \rho(A,k))^g$.
%
% Repeating this process, we then obtain $m \le g$ positive integers $M_i \le c_i \rho(A,k)^g (\log \rho(A,k))^g$, with each $c_i = c_i(g) > 0$, such that $[M_m\cdots M_1]P =0$. Set $c_0=\prod_{i=1}^m c_i$, which depends only on $g$. Then $M_m\cdots M_1 \le c_0 \rho(A,k)^{gm} (\log \rho(A,k))^{gm} \le c_0 \rho(A,k)^{g^2} (\log \rho(A,k))^{g^2}$. We are done.
\end{proof}

Finally, we drop the principally polarized hypothesis.
\begin{proposition}\label{PropDavidBound}
  For each integer $g\ge 1$ there is a constant $c(g)>0$ with the
  following property. Let $(A,\cL)$ be a
  polarized abelian variety of dimension $g$ defined over $k$. If $x\in
  A(k)$ has finite order, then
  $$\ord{x}\le c(g)
  (\deg_\cL A)^{1+(2g+1)\lambda(g)} \rho(A,k)^{\lambda(g)}.$$
\end{proposition}
\begin{proof}
  Let $k',B,$ and $\varphi\colon A_{k'}\rightarrow B$ be as in
  Lemma~\ref{lem:ppreduction} applied to $(A,\cL)$.
  Say $x\in A(k')$ has finite order. Then so does $\varphi(x) \in B(k')$ and
  $\ord{x}\le \#(\ker\varphi)\cdot \ord{
    \varphi(x)} \le \frac{1}{g!}(\deg_\cL A) \cdot \ord{\varphi(x)}$.
  We apply Lemma~\ref{lem:DavidBound} to bound 
  $\ord{\varphi(x)}$ from above and conclude the proposition after a
  short calculation.
\end{proof}

The value of the exponent $\lambda(g)$ is irrelevant for the
method. It suffices to know that it is a function of $g$.

%The proof of Proposition~\ref{PropDavidBound} 
%One can also prove Proposition~\ref{PropDavidBound} by applying \cite[Corollaire~1.5]{DavidMinHaut} for simple principally polarized abelian varieties, Poincar\'e's Complete Reducibility Theorem, and Masser--W\"{u}stholz's Isogeny Theorem   \cite{MW:abelianisog} to reduce to the principally polarized case. This argument is shorter, but uses % The proof we presented here is more elementary and does not use 
%Masser and W\"{u}stholz's deep theorem.

%%% Local Variables:
%%% TeX-master: "main"
%%% End:




\bibliographystyle{alpha}
\bibliography{literature}


\end{document}
